\section{Neural Signals and Physical Skills}

A hand can follow the wrong command perfectly. That difference between an issued command and an intended action is where NeuroTranslation begins: relate an identified neural recording to a specified body quantity, then check what happened. Predicting movement does not by itself identify the mechanism that produced it, and executing a command does not establish learning. The destination is faster physical learning that remains after teaching stops. The present evidence is bounded recording-specific software, separately calibrated simulations and a completed neural-decoding comparison whose improvement requirements failed. Decoding here means predicting a specified target from recorded neural input.

The LINK archive comparison makes the problem concrete. Two history-reading models lower average development error, but both fail the requirements that prohibit a harmed recording-day/finger-group and require positive neural information. All five separately defined benefit claims fail too. Source and scalar-calculation checks support that negative report; they do not qualify either model for use. The historical reaching, muscle-signal, attempted-finger and simulated-keyboard studies answer other questions. They are separate accounts, not replications that can repair LINK's failed comparisons.

\emph{NeuroTranslation: A New Paradigm In Brain-Body Mapping.} names the proposed approach, not an established priority claim or a new biological mechanism. Its models and simulators are not newly invented by being used here. \emph{NeuroTranslation} is the proposed field; \emph{NeuroTranslator} is the implemented software joining particular recordings, prediction models, playback and computer-body checks. The \emph{NeuroTranslation Network} is the proposed architecture for connecting more components while keeping each value's origin and meaning clear. \emph{Brain2Body}, with the numeral meaning ``to'', is the wider programme. Its companion article introduces the biological-interface proposal independently; neither manuscript uses a presentation or demonstration as scientific authority.\citep{brain2bodyworking}

Each study keeps its historical identity and terminology. The \hyperref[programme-definitions]{definition supplement} identifies this edition's shared registry version and fingerprint; a changed public label cannot retrospectively change the registry an earlier study used. The following comparisons ask what each recording can support, which outcome checks it, and where the connection fails.

\section{Measure Brain--Body Relationships}

The programme's starting point is that neural structure constrains function
without determining it. A connection map describes possible paths. It does not
supply every interaction strength, cell state, delay or learning rule. FlyWire's
adult female fly reconstruction is anatomical evidence with stated coverage and
reconstruction limits, not a complete behaving brain and body.\citep{flywire2024}

An electrode observes a system through an instrument. Scalp brain, spinal surface, peripheral nerve and muscle electrical recordings have different source meanings: EEG, ESG, ENG and EMG respectively. EMG records muscle activity, not peripheral nerve activity directly. Functional magnetic-resonance imaging, or fMRI, observes an indirect blood-oxygen-related signal. In the classic simultaneous-measurement study, its relation to local electrical activity depended on the measurement and site. It is not a photograph of thought.\citep{logothetis2001}

Name the source and target of every mapping, along with the subject, task, recording method, time support, units and tested domain. An electrode's contact identifier belongs to a declared source inventory; it does not establish that the same neurons were isolated on two dates. A neural feature is a specified numerical summary of a recorded signal.

Physical size, organizational level, computational detail and statistical averaging answer different questions. The 192 LINK input features are two measurements from 96 physical contacts, not 192 verified cells or 192 separate anatomical sites. Averaging four scoring cells equally defines a scoring rule, not a physical scale.

Perception is subjective. Movement is independently measurable through specified
positions, trajectories, timing, contacts and forces, each with calibration and
measurement limits. Sensory input guides action; action changes the body and
environment; feedback informs the next action. Feedback-control theory provides
a framework for testing that loop under uncertainty.\citep{todorov2002}
The present offline neural comparisons and simulated feedback loops remain
distinct from a measured living loop. A biological causal claim needs appropriate
intervention and outcome evidence in addition to predictive association.

Context also limits language-like reconstruction. The published semantic reconstruction depended on the participant and cooperation; it did not establish unique universal access to subjective content.\citep{tang2023}

The quantum-measurement analogy concerns observation and interaction. Quantum backaction means that measurement changes the state of the quantum system being measured. Quarks are subatomic particles. Neither backaction nor quark-scale measurement is offered here as an explanation of neural decoding limits at much larger physical scales.\citep{quantum2012} This programme claims no one-to-one decoding of private subjective experience.

Accuracy means agreement with an independent reference; precision means
agreement across repeated measurements under specified conditions. A repeatable
wrong output can be precise and inaccurate. A reference is not assumed
error-free: its acquisition, sampling, calibration and uncertainty remain part
of the test.\citep{nist1297} A correct guess is not a validated model.
A disclosed approximation may earn useful evidence in its tested domain.
Invented biological content presented as faithful reproduction is hallucination.
``Biological virtualization'' would require a digital model independently shown
to reproduce specified behaviour and causal dynamics within stated error bounds.
Neither a realistic surface nor a high simulator reward establishes that claim.

\section{Recording, Prediction and Skill Transfer}

A map is useful only if its inputs can be observed, its outputs can be checked,
and its meaning survives the conditions under which it will be used. We therefore
ask three separate questions. Can we observe useful signals? Can we connect them
to a particular body and a measured task? Does the connection remain useful as
recording conditions, bodies and tasks change? Each question requires its own
measurement. A clear drawing or a stable replay file cannot answer all three.

The software reported in the original evidence snapshot addresses selected
parts of those questions. It uses available macaque and human recordings, a
human arm simulation and a separate macaque arm reference. It does not establish
a lasting implanted interface or teach a living nervous system a skill.

A future experiment should begin with ordinary behavior and a defined baseline.
If a human-related biological feature or a tissue connection is studied, the
feature and comparison must be specified rather than treating a broad label as
an explanation. Experienced partners, computational trainers and disease-specific
conditions are possible future comparison families. They are proposals, not
outcomes of this computational work.

Task difficulty and assistance need controls. Simple straight routes and routes
with changing headings should be compared separately. Useful competitors include
movement-only predictors, signals shuffled in time, matched assistance and sensing
delay, unfamiliar sessions and animals, and performance after trainer instructions
are withdrawn. A convincing animation cannot replace those comparisons.

If a trainee needs an assistive bridge to express movement, that bridge should
remain standardized while trainer instructions are removed. Switching off the whole
bridge may prevent movement regardless of what the trainee retained. Assistance,
retained learning and tissue repair therefore need separate outcome measures.
A signal that remains useful in a saved file does not demonstrate that a living
recording interface remains useful over time.

For a proposed biological trainer, compare an experienced partner, a novice partner,
individual practice and a matched no-link condition, then test the trainee alone.
Replay comparisons should distinguish correctly aligned signals from stale or
unmatched signals and an inactive comparison. Unfamiliar tasks and bodies must
stay outside fitting. A new-looking maze or an equal target score does not by
itself establish equal difficulty or retained learning.

Body observation also needs independent checks. For the proposed hand--arm task,
synchronize the left and right tripod views and compare what each reveals or
misses. A separate perpendicular depth-sensing view could be compared for
coverage. This software did not acquire or calibrate either arrangement.
Transparent layers and cutaways explain a model; they do not turn hidden nerve
activity into a measurement.

\section{Recording-Specific Models}

The earlier computational studies explain how the recording-specific software
works and where it fails. The current LINK cohort and its completed negative
comparison follow. Each older study keeps its original tested scope; support
already inspected cannot become fresh confirmation.


The earlier software snapshot had seven executable translation routes. Each route belongs to an
identified recording and a particular signal quantity. A packet carries recent
values, their recording identity, channel order, quantity, units, body frame
where relevant, and optional sample times. The translator returns a named
estimate with an explicit time and source. It never silently inserts missing
spinal or peripheral motor signals.

\begin{figure}[t]
\centering\makebox[\textwidth][c]{\includegraphics[width=160mm]{figures/architecture.pdf}}
\caption{Historical diagram: the implemented chain and the proposed next stage. Independent
recordings keep separate identities and clocks. Body feedback checks simulated
movement; it does not reveal hidden nerve activity.}
\end{figure}

One route estimates fingertip velocity from a macaque movement-brain recording.
A second estimates 16 normalized muscle-activity channels from a different
macaque recording. A human route classifies attempted finger identity from
movement and parietal brain areas. Three touch-brain routes estimate earlier
mechanics or expected brain activity from mechanics or joint history. The seventh
estimates velocity on a genuinely different macaque day after separate calibration.
These routes share a software interface. They are not one simultaneously
observed brain, spinal cord, nerves and body.

The seven routes give the software a common way to handle different recordings. Their separation is essential: combining the interface cannot supply simultaneous observations that the recordings lack. This recording-specific implementation, defined in ``Neural Signals and Physical Skills'', establishes neither a universal trained decoder nor the proposed larger network or biological teaching. Older technical names remain glossary aliases.

The scale distinctions in ``Measure Brain--Body Relationships'' matter here. A small prediction model does not imply a small measured biological structure. Physical size, organizational level, computational detail and statistical averaging remain different axes; averaging output groups cannot supply an anatomical measurement.

\section{The current LINK development cohort}

What does recent neural history add when every model must answer the same position question? The LINK study uses DANDI 001201, version 0.251023.2336, Monkey N in the
centre-out finger task: one rhesus macaque, six recording sessions, 375 trials
per session.\citep{linkarchive} Four dates are training data (January 27 and
30, February 4 and 5, 2020); two later dates are development validation
(February 6 and 11). Two further February dates and two June 2023 drift
recordings remain reserved and closed. These are distinct future confirmation
supports; they are not part of the reported development results.

Ninety-six physical contacts in hand primary motor cortex provide 96
threshold-crossing features and 96 spiking-band-power features. Threshold crossing
counts detected events crossing a chosen threshold; spiking-band power summarizes
signal power in the source's selected high-frequency band. The counts use
$\log(1+x)$, a transformation that compresses large counts while retaining zero.
The archived power features keep their supplied scale. Absolute calibration is
unverified, and an archive voltage label conflicts with the publisher's
threshold-count description. File labels do not resolve that conflict.

The outputs are two archived normalized sensor positions: index finger and the
combined middle, ring and small-finger group. We call the latter the combined
finger group rather than implying three independently measured trajectories.
These targets are not radians, newtons, individual joint angles, muscle force,
EMG, peripheral nerve activity or spinal activity. Conversion to a body model
requires separate physical calibration and an explicitly tested interface.

That comparison starts before fitting a predictor: even the input scale must come from training alone. The scaler's input means and population standard deviations were fitted to 107,605 unique trial-owned training frames. Population standard deviation measures spread around the mean using the count of all fitting values as the divisor. Validation supplied no scaler statistics.

Every candidate uses the same eleven-frame history inside one trial, with the following neural frame added where the auxiliary target requires it. The position label belongs to the final input timestamp. This gives 91,105 training endpoints and 45,757 development-validation endpoints. The eleven timestamps span about 199--201 milliseconds. All six series in each file share the recorded clock, but unresolved bin anchoring and filtering prevent that span from being treated as physiological input support or verified online delay.

\begin{table}[t]
\centering\small
\begin{tabular}{lrrrr}
\toprule Date & Role & Trials & Source frames & Eligible endpoints\\\midrule
2020-01-27 & Train &375&27,649&23,524\\
2020-01-30 & Train &375&26,944&22,819\\
2020-02-04 & Train &375&26,441&22,316\\
2020-02-05 & Train &375&26,571&22,446\\
2020-02-06 & Development &375&26,957&22,832\\
2020-02-11 & Development &375&27,050&22,925\\\bottomrule
\end{tabular}
\caption{The current one-animal cohort. Each session has 4,125 frames that are
not prediction endpoints but can supply declared history, auxiliary labels or
training scaler statistics. Endpoint counts enumerate archive support; they
are not independent subjects, independent windows or confidence intervals.}
\end{table}

A saved independent source audit reconstructed the six sources and the prepared
packet values. Its scaler difference was at most $4.54\times10^{-13}$ within the
frozen criterion. That evidence authenticates the named preparation. It does
not demonstrate accurate decoding, causal physiology or independence of the
common learning engine. The separate human EEG/\allowbreak{}ESG/\allowbreak{}ENG/\allowbreak{}EMG replay is an evoked
stimulation source on its own clock; it is not input driving this macaque study
or the historical macaque-to-human-arm simulation.


\section{Recorded Signals and Missing Measurements}

The earlier sources extend the range of targets, not the measurements available in LINK. The public reaching archive pairs macaque movement-brain signals with fingertip
position. The original task release supplies a processed fingertip-velocity target;
the independently acquired day supplies recorded positions from which we
derive fingertip velocity. The touch-brain study adds hand motion, handle force and turning
force, with joint angles and muscle lengths obtained through tracking and body
analysis. Those modeled lengths are not measured living muscle fibers.
\citep{odoherty2020,chowdhury2020,nlb2021,falcon2024}

The human finger recording contains cues and attempted finger labels. It does
not provide executed finger trajectories, sewing, cube manipulation or measured
muscle forces. A separate human source contains simultaneous scalp brain,
surface spinal, peripheral nerve and muscle signals during nerve stimulation.
Its companion lacks muscle channels. These surface signals are useful paired
observations, but they are not individual spinal-cell spike trains or evidence
for voluntary whole-body control.\citep{human_fingers2023,spinal_dataset2023,spinal_companion2023}

The historical curated catalog had 37 resource records. A catalog entry does not mean that
we acquired every file or trained on it. It includes functional recordings,
movement data and anatomical connection maps with different licenses. An
anatomical map describes connections; it cannot supply unrecorded activity over
time. Sources remain identified by release, asset, bytes and checksum.

The historical second reaching day was chosen before this study inspected its raw data or decoder results: the closest eligible published Indy day within the initial download limit, excluding the existing day. Its 208,792,806-byte original file was checked against the publisher's fingerprint and the stronger local fingerprint. It contains 171 nonempty sorted signal features; the unsorted residual channel is excluded. Channel numbers do not establish that the same cells were recorded across days. Because the selected day precedes the original model-training day, this is a look-back comparison, not a future-day test.

We divide the recording into short completed intervals, called bins.
Every completed 20-millisecond bin uses its last available 250-times-per-second
fingertip sample. We subtract successive positions and divide by the elapsed
time to obtain velocity. The actual
position sample ends four milliseconds before the neural bin ends. That offset
comes from sampling; it is not a measured nerve-conduction delay. Cursor and
fingertip positions can differ by 11.9 millimeters and remain distinct.

\section{Compare Against Simple Models}

If the history helps, its advantage should survive comparison with a single sample and with a weighted sum of the same history. The frozen LINK comparison includes a constant training-target mean, a linear
model using the last sample, a linear model using the eleven-sample history,
and two matched gated recurrent-unit models. A linear model computes weighted
sums; ridge fitting adds a penalty on large weights. Its five predeclared
penalties are 0.01, 0.1, 1, 10 and 100, with an unpenalized constant term. Both
linear models select 100 using the same two development dates.

A gated recurrent unit is an established sequence model that carries a numerical state from one input step to the next; learned gates control how that state is retained or replaced. This is memory in the software, not evidence of biological memory. Here it has 192 inputs, one memory layer with 128
values, a two-position output and an additional next-neural output. The two
candidates retain the same full 169,090 parameters, although only 144,322 are
active in the position-only objective. The auxiliary candidate adds 0.1 times
next-neural mean squared error to position mean squared error. An objective is
the numerical loss minimized by software fitting. It is not a learning rule
shown to operate in the animal.

Both candidates use three fixed seeds, 20261007, 20261008 and 20261009, with matched initialization and sample order. Every run retains all 40 epochs, each a pass through the supported training windows. Batch size is 256, and the final partial batch is retained.

The established Adam update uses learning rate 0.001, averaging constants 0.9 and 0.999, numerical stabilizer $10^{-8}$ and no weight decay. Weight decay pushes model weights toward zero. Dropout randomly omits internal model units during training. A gradient measures how a small weight change affects the loss; gradient clipping limits the size of the update-direction vector. Target clipping replaces out-of-range targets with chosen bounds. None of these four controls is applied.

The reported study configuration specifies four mathematical worker threads and one framework scheduling thread.

Each seed selects its lowest equal-date/group validation error, taking the
first epoch on an exact tie. The final predictions average all three selected
seed models equally; a favourable seed is never substituted for that ensemble.
The measured input condition and matched association and temporal-order
controls are distinct. Association controls break the relevant pairing;
order controls challenge use of input sequence order. Neither control enters
the best-model selection contest. Two candidates, three seeds and three
conditions give 18 runs and 720 saved epochs.

The primary error averages mean squared error for two groups on two dates: four equally weighted scoring cells. Mean squared error averages squared prediction-minus-reference differences, in squared source-normalized position units.

The whole-trial paired bootstrap repeats the error comparison 1,000 times. It resamples complete trials within each date and shares the draws across models and groups, preserving correlated frames inside a sampled trial. A percentile is a cut point below which a stated fraction of values lies. The resulting percentile interval is a descriptive 95-percent interval conditional on these two dates and the validation-selected models. It does not confirm performance on a new animal or recording.

A development improvement must meet \emph{all} original conditions: lower
aggregate error than the selected ridge; no higher error in any date/group;
a negative upper interval bound for candidate-minus-ridge error; and positive,
defined ordinary $R^2$ in every cell that also exceeds its matched association
control. $R^2$ compares squared errors with a prediction of that evaluation
cell's observed mean. The separately reported training-mean skill score uses
a different constant reference and is not interchangeable with ordinary $R^2$.
History, order-use and auxiliary-benefit claims retain their own conjunctions.
An improved average cannot override a failed cell or rewrite those criteria.


\section{Why the earlier implementation used small models}

A larger model would change the cost of these comparisons; it would still have to pass them. The earlier environment receipt identified a Ryzen AI Max+395 processor and roughly 125 gibibytes of memory. Those training runs used four processor threads; graphics-processor training was unverified, and no paid cloud training was used. The current LINK study retains its specified CPU route. Its reported bounded GPU prefix took 4.565 seconds end to end, versus 2.427 seconds on CPU, and remains unpromoted. That prefix qualifies neither full training nor energy or cost savings. Compact models made training runs, controls and replay possible in that environment.

The first movement model has 232,388 adjustable numbers; the muscle-activity
model has 195,216. Both first learn to predict the next signal bin from recent
signal history, then learn the paired movement or muscle target. This is real
training followed by fine-tuning. It is not a claim that a public pretrained
brain foundation model was downloaded and adapted. The new-recording model
trained from scratch has 256,045 adjustable numbers.

Large-scale neural models already exist. Some read spikes, others scalp
electrical signals or magnetic-resonance measurements. Their inputs differ,
and a larger model does not make those measurements interchangeable.
\citep{ndt2_2023,poyo2023,brainlm2023,neurogpt2024}
The question for a larger shared model is whether it beats well-controlled
specialists on unfamiliar recordings and bodies. Our first shared movement
and muscle candidate failed its validation comparison and remains rejected.

\section{How the earlier predictions were evaluated}

Validation chooses the model; testing checks the result of that choice. Training examples adjust its parameters. Validation examples choose between candidates and decide when training stops. Test examples describe performance after that choice. The continuous movement split is 70/15/15 percent over time,
with a one-second gap at boundaries and complete signal histories. The muscle
and touch-brain splits keep whole trials together. Our normalization is fitted
only on training examples. Source-provided filtering or normalization may
already have used later data, so our split does not certify the entire upstream
recording pipeline as causal.

The prediction score $R^2$ compares a model with predicting the evaluation-set mean. One means perfect agreement; zero matches that reference; a negative value is worse. Scores are averaged across named output channels or defined groups. They are neither percentages of learned skills nor official benchmark rankings.

Shuffled inputs and simple linear predictors supply additional comparisons. A linear predictor is a fitted weighted sum of input values. ``Typical error'' below squares the errors, averages those squares and takes the square root. In the two-joint tracking case, each frame combines both joint-angle differences before that summary is calculated. Radians measure angle; meters and millimeters measure length.

The touch-brain reconstruction compared zero, 40 and 80 milliseconds of
alignment on common usable examples. Validation selected 40 milliseconds:
brain activity reconstructs \emph{earlier} mechanics. It is not a future movement
command or proof of biological conduction time. Its separate reverse model
predicts positive expected count means. Those means are statistical estimates,
not new recorded spikes or stimulation instructions.

The different-day study fixes two supervision amounts before outcomes are
known. The small amount uses 1,425 new training labels, or 28.5 seconds, \emph{plus}
6,000 validation labels, or 120 seconds. All methods can use 570.14 seconds of
new training-brain signal history for signal-only learning. The original model
also carries 22,598 paired training and 4,723 validation examples from its
earlier study. Thus ``first 5 percent'' is not five percent of all supervision
or evidence of biological mastery in 28.5 seconds.

Four candidates receive the same new paired examples in each supervision
amount: an input adapter into the old frozen model, the same-sized adapter into
a randomly initialized frozen model, a newly trained model and a simple linear
predictor. Frozen transfer is accepted only if its validation score exceeds
every control by at least 0.01. The full amount is the interface default;
both evaluated amounts were fixed before outcomes. All eight outcomes share
the same 6,050 local test windows, now inspected. They are not eight independent
prospective experiments.

\section{LINK Results: Neither Model Passed}

Across 18 runs and 720 saved epochs, lower mean error was not enough: neither candidate passes the frozen improvement conjunction. The postfit audit checks the named sources and scalar calculations in this negative report. Postfit means checking reported outputs after fitting; it does not independently certify final-checkpoint or full optimizer arithmetic.

\begin{table}[t]
\centering\small
\begin{tabular}{p{.40\linewidth}p{.28\linewidth}p{.24\linewidth}}
\toprule Model & Development mean squared error & Outcome\\\midrule
Training-target mean &0.0317577871&Reference\\
Last-sample ridge &0.0314982382&Selected reference\\
Eleven-sample ridge &0.0383765124&No history advantage\\
Position recurrent ensemble &0.0312304247&FAIL\\
Position plus next-neural ensemble &0.0299609175&FAIL\\\bottomrule
\end{tabular}
\caption{Saved reported scalar summaries from one macaque and two development
dates. Lower aggregate error alone is insufficient. Both recurrent candidates
fail the complete original development gate, and no model is promoted. These
figures were transcribed from accepted source-bound reports, not recalculated.}
\end{table}

For the position-only ensemble, the difference from ridge is
$-0.0002678134$, with conditional 95-percent interval
$[-0.0014148839,\,0.0009468841]$. For the auxiliary ensemble, the difference is
$-0.0015373207$, with interval $[-0.0027586737,\,-0.0003470099]$. Its aggregate
interval excludes zero. Its no-harm and neural-information conditions still
fail. The aggregate result and the failed conditions must be read together.

On February 11, position-only error exceeds ridge by 0.0052471986 for index and
0.0061405635 for the combined group. Auxiliary error exceeds ridge by
0.0043574767 for index. The position-only ordinary $R^2$ values on that day
are approximately $-0.099742$ and $-0.353799$; auxiliary values are
$-0.071162$ and $-0.116390$. The corresponding association-control values
stay visible in the ledger. Exceeding a negative control in one comparison
cannot satisfy a requirement for positive information in every cell.

All five additional claims fail: history benefit for position-only, order-use
for position-only, history benefit for auxiliary, order-use for auxiliary,
and auxiliary-objective benefit over position-only. The original selected
linear models also retain negative ordinary $R^2$ cells. Neither averaged
performance nor one agreeable trajectory establishes reliable coverage of all
evaluated support. These are development observations, not population estimates,
cross-species translation or biological teaching.

The auxiliary ensemble exposes what the average concealed. Its aggregate interval excludes zero, yet its average improvement coexists with a harmed output and negative information scores. Under the unchanged requirements, that is a failed development comparison. The four reserved recordings remain closed, so this result supplies no confirmation beyond the development support.

The negative result leaves a separate implementation question open: were the full learning calculations independently verified? Both routes share Torch automatic differentiation, which computes changes in the objective, and Adam, which uses those changes to update fitted state. Source review and saved numerical checks of the report do not independently verify that common learning engine. An independent numerical route must still check the full learning calculation: complete NumPy calculations, restored random-generator states, reload equivalence, final-checkpoint forward calculations, derivatives, Adam states and complete update trajectories. NumPy is a numerical-array library. Those checks remain outside admitted scope; the final-checkpoint/update gate and current fitting remain held. No model or GPU route is promoted. The documentary details are retained in ``Frozen LINK Development Methods''.

\section{Failure Analysis}

Where did the auxiliary ensemble lose its advantage? The saved residual diagnosis decomposes the auxiliary model's second-day index
excess mean squared error. About 78 percent is the difference in squared mean
error and about 22 percent is the difference in centred residual variance.
Its mean prediction is 0.058645 below the target in the source-normalized
sensor units. This describes the miss; it does not establish a cause. Trial
composition, state-dependent offsets and recording/calibration changes remain
competing explanations. Centring predictions on validation answers would fit
those answers and invalidate the intended comparison.

A separate training-only context analysis uses the mean of ten preceding neural
bins in each existing eleven-bin window. Across four dates, 1,500 trials and
91,105 supported endpoints, summed threshold-feature variance is partitioned
into 71.56 percent within trial, 18.85 percent between trials within date and
9.59 percent between dates. The corresponding power-feature shares are
51.28, 13.55 and 35.17 percent. These ratios weight features by variance and
are descriptive archive summaries, not uncertainty intervals or proof of
biological state or neural drift. Equal date/trial/frame-within-trial weighting
gives similar fractions. The existing recurrent model already sees these bins;
variation alone cannot demonstrate the usefulness of an added branch.

Joining the exact source indices, times and trials leaves 89,605 endpoints with one supported prior sensor sample and 88,105 with two. Trial beginnings without labels remain missing; no labels are borrowed across trials or interpolated. Each training date supports all sixteen joint prior-position quartile cells. A quartile is one of four divisions of an ordered numerical distribution.

Correlated frames are not independent samples. Body-position strata are groups of body positions used as diagnostic labels. They can locate a neural-only predictor's failures; supplying those labels to the predictor would change its input condition, so they cannot enter it.

The first descriptive calculation failed two skewness checks at the unchanged
absolute and relative $10^{-12}$ tolerance. Skewness measures asymmetry of a
numerical distribution. The corrected second version changes only the
per-date accumulation of that quantity; 768 values pass the named independent
checks and all 119,848 other numerical leaves are unchanged. The rejected
version remains in the evidence record. Acceptance concerns these descriptive
calculations, not model improvement or the final-checkpoint gate.

The proposed next readout must obtain its context from permitted neural history. It would let past neural context change the decoder's output through 780 additional parameters. Its rank-four branch uses four learned context summaries to modulate that output. Rank describes the restricted numerical representation, not anatomical scale. The branch needs constant-context, association and order controls.

Selection must stay strictly inside training, using independently verified whole-trial inner folds: training trials divided into internal fitting and checking groups. Scaler statistics must come from unique inner-fit bins. Renormalizing rounded packets that were already normalized across training cannot establish that independence.

The branch is unexecuted, and current fitting is unadmitted. A prospective test must preserve initialization, random seeds, selection, uncertainty, complete interface/derivative/update evidence and the four closed reserved dates. Body feedback, reward-driven learning and broader extensions remain proposals, not LINK outcomes.


\section{Earlier prediction and transfer results}

The earlier studies answer narrower prediction questions, with different targets and recordings. The original movement specialist scores 0.729 against 0.559 for its linear
control. The corrected 16-channel muscle specialist scores 0.704 against
0.595; its weakest channel scores about 0.088. Source muscle activity has
already been normalized. A stale source file label saying millivolts does not
turn those dimensionless values into physical amplitude.

The human selected classifier correctly identifies 45 of 50 attempted-finger
trials, or 90 percent; its recent-signal model scores 68 percent. That evaluation
run was inspected in an earlier pilot and is development evidence. A separate
uncalibrated day comparison scores 10 percent, matching ten-class chance.
It does not establish executed finger control or cross-day skill teaching.

The touch-brain decoder scores 0.710 over five named mechanics groups. The
reverse mechanics-to-brain model gains 0.177 bits per spike over a training-mean
reference, and the seven-joint-only version gains 0.119. This information score
compares how well expected counts explain recorded counts. It does not prove
that the model can generate a complete biological spike process.

\begin{figure}[t]
\centering\makebox[\textwidth][c]{\includegraphics[width=160mm]{figures/recorded_replay.pdf}}
\caption{Historical study: actual different-day brain counts and a fingertip-velocity component
predicted by the selected new model. The reference component is calculated from recorded fingertip positions.
The rows belong to this recording; they are not matched cells across days.}
\end{figure}

\begin{table}[t]
\centering\begin{tabular}{llrr}
\toprule New examples & Method & Selection score & Local test score \\
\midrule
First 5\% & Use the old fixed model & 0.393 & 0.425 \\
First 5\% & Use a random fixed model & 0.312 & 0.331 \\
First 5\% & Train a new model & 0.533 & 0.520 \\
First 5\% & Simple fitted formula & 0.276 & 0.220 \\
Full & Use the old fixed model & 0.683 & 0.704 \\
Full & Use a random fixed model & 0.631 & 0.646 \\
Full & Train a new model & 0.724 & 0.747 \\
Full & Simple fitted formula & 0.550 & 0.548 \\
\bottomrule\end{tabular}

\caption{Historical study: all eight different-day outcomes. The prediction score is defined in
the text. Both supervision amounts also require validation examples, recent
brain-signal history and, for the transferred model, prior training. The
local test examples are shared across methods.}
\end{table}

The old model carries useful information compared with a random frozen model, but training a new model wins at both supervision amounts. Both frozen-transfer candidates remain rejected. After local calibration, the selected full-amount new model scores 0.747. That result supports a recording-specific specialist, not zero-shot transfer or shared neuron identity. Different recordings and velocity preprocessing prevent treating it as an improvement over the original 0.729 case.

\begin{figure}[t]
\centering\makebox[\textwidth][c]{\includegraphics[width=160mm]{figures/transfer_comparison.pdf}}
\caption{Historical study: transferred information beats its random control but does not beat
the newly trained model. These are point estimates from one fixed experiment,
not uncertainty intervals across repeated random seeds.}
\end{figure}

\section{Separate human intention and simulated-hand studies}

The hand studies let us separate two questions that a successful animation can blur: did the neural decision identify the attempted finger, and did the body execute the command it received? These October 4 studies are separate from LINK and from the earlier reaching tests. Their saved acceptance scopes do not change the later LINK failures.

\subsection{Recorded brain signals select an attempted finger}

The later study used public recordings from one participant with tetraplegia,
a condition affecting movement in all four limbs. Its labels identify the finger
the participant was cued to attempt. They are not measured finger movements.
The three held recording dates were May 16, May 20 and May 23, 2016. Those dates
contain posterior parietal cortex recordings, from a named brain area. They do
not supply paired motor-cortex, spinal, peripheral-nerve or muscle recordings.

The model used eighteen completed 50-millisecond bins, covering 0.1 to 1.0 seconds
after the source Go signal. A decision could first be released one second after
Go. Fifty early labeled examples from each new date supplied local calibration:
examples used to adapt the model to that date. The same participant supplied all
three dates; this is not a study of 150 people or an uncalibrated new person.
Recorded unit identities remained local to each date. Missing observations were
kept unavailable rather than filled with invented measured zeros.

Seventeen fixed model candidates were compared. Selection used earlier validation
dates before the later held examples were opened. The selected simple classifier
correctly identified 116 of 150 later attempted-finger labels: 35, 44 and 37 of
50 on the three dates. A population model without those local labeled examples
correctly identified 11 of 150. These results do not show that a larger model is
better, that the same cells were matched across dates, or that the participant
learned a physical skill through replay.

\subsection{Neural Decision Errors and Muscle Responses}

The actual classifier decision, including errors and unsupported fingers, set
the body command. The system did not choose the best right-finger class after
seeing the answer. Four right fingers could issue goals to four physical keys;
left fingers and the right thumb remained unsupported. Reference labels scored
the learned replay and did not repair its predictions.

A saved learned policy supplied 39 muscle excitation commands to an openly
available, generic human-hand simulation. The wrist, thumb and finger abduction
were braced. The goals used fixed fingertip press targets and a 0.5 second pulse.
Body positions and velocities were initialized once per recording, then advanced
through physical integration without pose or velocity overwrites. Contacts, key
travel and release were measured in the simulator.

All 54 issued supported goals were physically executed. However, only 41 of the
60 supported right-finger source intentions produced the correct contact, and
13 unintended presses remained. The two fractions answer different questions:
54/54 measures whether the body followed the issued command; 41/60 measures
whether the recorded attempt led to the right physical result.

The fixed controls retained all source trials. Zero muscle excitation produced
0/60 source-intention contacts. A reference control given the true source labels
produced 60/60. Deliberately reassociating neural decisions with different trials
produced 4/60, despite executing all 54 issued supported goals. This reassociation
is an offline control, not a new chronological neural recording.

The probabilities were computed offline and released only after their original
source windows completed. The result is paced recorded replay, not live brain
acquisition. The participant did not perform the simulated motion or play music.
These results do not establish accurate hidden nerve activity, a learned piano
skill, biological teaching or a universal brain--body map.

\subsection{A separate simulated hand learned phrases and chords}

That replay asks whether a recorded decision can drive a press. The separate phrase study asks whether a goal-driven muscle policy can complete a requested sequence. It trained on requested musical phrases in a generic four-key simulated hand with braced wrist, thumb and finger abduction (sideways spreading), driven by scripted smooth fingertip goals. Recorded neural clips did not drive this policy.

The 32 episodes combine eight phrase orders, two tempos and two starting poses; they are not 32 subjects. The policy received requested notes, timing and fixed smooth fingertip goals; it received no teacher actions during held evaluation. New phrase orders, rhythms, two-finger chords, tempos and starting poses were reserved before tuning.

The selected corrected policy completed all 32 held phrases: 704 intended
strikes, 192 chords and timely releases, with no extra presses. The starting
policy completed 0/32, the first new fit completed 20/32, the named reference
controller completed 24/32 and zero excitation completed 0/32. Independent
numerical replay verified the saved results. These are constrained simulation
results. They do not show that a living hand learned music or that brain signals
encoded those phrases.

The original evaluator exited with code 143 after the closed results and complete final output had been saved. The termination's cause remains unknown. The independent full replay verifier exited 0; that numerical check does not make the simulator a biological reference. Spring-key limits are soft numerical constraints, so travel can exceed the nominal 6 millimetres. Saved poses and travel were neither clamped nor fabricated for display.

Local training of that muscle-policy study took about 2 minutes 36 seconds.
The later neural candidate training and validation took about 24.7 seconds.
Those are measured software runtimes; acquiring recordings, supplying labeled
examples and verifying physics require additional time. No completion time for
biological teaching can be inferred from them.



\section{The earlier muscle-driven arm experiments}

The arm study tests another link: can decoded movement goals produce simulated muscle-driven motion? The working human arm uses 63 simulated muscles that produce force. Their
activation changes over time rather than jumping instantly to the command,
with 500 simulation steps per second. Two predicted fingertip-velocity components are explicitly
mapped to shoulder and elbow movement goals; other independent joints are
braced. The simulator calculates each next body state from muscle forces and
the current state. Joint poses are not overwritten, and extra artificial
joint-turning forces do not replace the muscles.
\citep{mujoco2012,myosuite2022,myosim2026}

Feedback correction reads simulated body state. Recorded-position-derived
movement is reserved for scoring. In the bounded 500-frame new-day replay,
typical corrected tracking error is 0.0374 radians against 1.5803 for raw
control. This comparison follows each branch's own command target, which can
change under the correction's bounds. Against a common scoring reference
mapped from recorded positions, the errors are 0.26302 and 1.49325 radians.
These are numerical body-control results, not neural accuracy or measured
macaque muscle behavior.

The corrected arm's largest movement beyond a declared joint limit is 0.000406 radians,
below the declared 0.001-radian numerical tolerance. Raw replay violates that
tolerance in 75 of 500 frames. Passing the tolerance means passing this solver
check. It does not establish biological safety or physiological fidelity.

The separate macaque model has seven coordinates, 39 muscle paths and a rigid hand. Its retrieved source revision postdates the touch-brain paper, so it belongs to the same source family rather than being a certified copy of that study's specimen model. The original muscle-rate parameters contain zeros. A separately named copy changes 117 muscle settings to the software's official defaults. That documented assumption does not establish living-muscle behaviour.\citep{opensim2018,limblabmodel2026}

Three original-model muscle controllers remain rejected. The force-feedback
case has per-coordinate typical error 0.05505 radians versus 0.04512 for a
constant muscle command. This error squares all joint-angle errors, averages
them across coordinates and time, then takes the square root.
An initial-speed protocol correction is retained, and this already inspected
segment is not called a fresh test. The later controller plans short movements
in isolated copies of the simulated body state. It loses 0.05971 to 0.03263
radians despite improvement on constructed simulation targets. Planning
takes about 4.52 seconds for a 20-millisecond interval and is offline.
Its cached replay can play at normal speed; that is not real-time planning.

\section{Anatomical Templates and Participant Registration}

Before a neural overlay can identify a participant's anatomy, its template must be checked and registered to that participant. The retained reconstruction audit reproduced 122,284 template vertices, 244,760 triangles and 81,924 atlas values, with three source MRI planes. Vertices are mesh points, triangles connect them into a surface, and the atlas assigns template regions. Both mid-cortical surfaces passed the specified closed-surface checks, including Euler characteristic $+2$, a connectivity summary appropriate to those tested surfaces.

In that audit, the optional region boundary had 432 open edges, 838 winding inconsistencies---inconsistent triangle orientations---and six invalid vertex fans, defects in neighbourhood connectivity. A second reader using the established nibabel imaging library was unavailable. Successful counts elsewhere do not repair these defects.

Those checks establish named template/source correspondence. The recorded subject still needs independent spatial registration; the checks neither localize individual neurons nor demonstrate physiological function. Registration means independently establishing spatial correspondence
with the actual subject. A group template may contextualize ``hand motor
cortex'' while an electrode's local identity remains local. A region colour
cannot claim that an absent peripheral pathway was measured.

A template cannot fill those measurement gaps by becoming more lifelike. Every visual body and neural map still needs a lifelike, tangible three-dimensional representation refined against a \emph{series} of identifiable real-world references. Those comparisons must cover recording site, time, pose, physical scale and anatomy, with independent review of both the parts and their assembly. Biological identification colours, cutaways and useful views can explain those relationships. Functional overlays must also preserve the displayed quantity, clock, availability and measured, predicted or simulated origin; otherwise the view changes what it appears to claim.

The required held-out trajectories, joint and force measurements and complete anatomical and functional validation remain missing. The four historical figures are evidence diagrams and plots. Retaining them preserves those earlier comparisons; it supplies neither newly graded 3D biological models nor a fresh visual review.


\section{Earlier source and body defects remain evidence}

The earlier source audit shows why a plausible score cannot authenticate its input. It found that 33 of 64 brain-event streams in the muscle study were
not globally sorted. Counting them as though they were sorted corrupted
866,594 entries in the table of channel-by-time counts. These are count-table
entries, not 866,594 living cells. Earlier muscle, shared-model and uncertainty
results were withdrawn.
The corrected historical streams were sorted, checked against independent whole-source counts
and retrained with the unchanged seeds and split rules. Initial green tests
had not proved source correctness.

The same demand for source checks reaches into the body geometry. A separate geometry audit found an implausible native triceps detour around
the elbow. A rule based on the original attachment locations restricts three
wrapping sections to the part of each muscle path between the upper-arm bone
and the ulna, one of the two forearm bones. It removes 14
off-surface contacts in seven inspected anomalous frames, without deleting
muscles, moving attachment points or fitting anatomy to those frames.

A global size factor is then fitted only on 128 training poses:
\begin{equation}
 a = \frac{\sum_{t,i} L_{t,i}^{\mathrm{model}}L_{t,i}^{\mathrm{reference}}}
              {\sum_{t,i}(L_{t,i}^{\mathrm{model}})^2}.
\end{equation}
Here $a$ is one common multiplier; $t$ identifies a training pose; $i$ identifies
one of the 39 muscle paths; and $L$ is its length in meters. The equation chooses
the multiplier that minimizes squared length error on those training poses.
The reference lengths were generated by the recording's tracking/body analysis,
not measured directly from living fibers.

\begin{figure}[t]
\centering\makebox[\textwidth][c]{\includegraphics[width=160mm]{figures/length_fit.pdf}}
\caption{Historical study: one size factor fitted on training poses reduces length error against tracking/model-derived references. These reported values remain those of the earlier study. Agreement with those reference lengths does not identify specimen anatomy or establish living muscle behavior.}
\end{figure}

The fitted factor is 1.2413815. Length error drops from about 27.75 to
0.552 millimeters on separate validation/test poses. A later, separately named
model applies that factor through official body-model scaling rather than
multiplying displayed paths after export. Its native lengths reproduce the
scalar predictions within 6.13 micrometers. The modeled masses are retained;
the geometric changes to inertia---resistance to a change in turning motion---
and muscle lengths remain assumptions.
Dynamics are blocked on that scaled viewing candidate.

The 400 displayed observed and decoded poses are reused diagnostics. Their
small observed-pose length error is not evidence for accurate decoded movement.
The known wrap anomaly also occurs later than the evaluated lookahead-controller
window, so it is not proved to explain that controller's rejection.

\section{Earlier interface checks and their bounded scope}

An interface can reject a malformed packet without knowing where its values physically came from. The interface checks recording identity, signal shape and quantity, channel
order and, for mechanics, names, units and body frame. When times are supplied,
it checks complete-bin cadence and the last sample time. Without them,
nominal cadence is assumed and reported. Integer spike counts are checked
before conversion to the model's lower-precision numbers. These rules validate
the caller's declared contract; they do not independently authenticate the
physical origin of supplied values.

The model outputs remain labeled as estimates. Expected counts cannot silently
become measured spikes. Missing motor-spinal and peripheral-motor observations
remain unavailable. Controller correction acts on simulated movement, not on
recorded nerve counts. The saved different-day cache contains actual body
positions, activations, geometry and muscle paths. It contains no saved joint
speed array; no such array is fabricated from differences in position.

In the historical verification, all seven routes also ran through exported NumPy weights. NumPy is a numerical-array library; this route calculated outputs without importing the main training library. A separate historical live verifier checked all seven examples within numerical tolerance 0.000002, preserving exact response identity, units, times and support descriptions. The checked route rejected mismatched requests. These checks establish replay and interface consistency, not a universal neural language.

The earlier dashboard snapshot showed actual recorded counts, estimates, model choices, body
state, activation and playback. Its natural-language task interface supported
explicit reaching lessons and measured reach-to-grasp replay. The practice
software example adjusted two simulated controller gains and displays the measured
simulation errors and rewards. Sewing, cube solving and breakdancing lack
supporting recordings and task physics here. They remain future composite tasks.

\section{Current Viewer and Missing Measurements}

Seek to another moment in a replay and the displayed quantity must still mean the same thing. In the retained cached-hand component check, the graphical user interface (GUI) loaded an identified saved replay with 336 frames at a 20-millisecond source interval. Play, pause, seek, stepping and restoration preserved that source. An uncached selection triggered no automatic replay or run request. Missing values appeared as dashes, and unavailable or failed sources cleared stale replay and geometry. These behaviours matter because an old hand pose beside a newly selected recording would describe a relationship the source never supplied.

The saved development example contains 22 of 22 strikes, 6 of 6 chords and 22 of 22 timely releases, with a 34-millisecond worst timing error and no extra strike. It is one constrained simulated example. Separate human neural clips did not drive that goal-driven hand policy. Its snapshot-level check does not establish current operation of every GUI route.

The original-control checks cover bounded source switching, gaps, replay controls and camera/restoration behaviour. Complete validation still requires the original source-time, geometry, materials, renderer-lifetime, functional-map, anatomical-reference, device and offline checks. No whole-GUI acceptance or performance improvement follows from the checked component.

Opening a package is a separate question from inspecting a replay. The online candidate needs its whole directory and assets; the original-source static offline export is a distinct artifact. Views that require the application's service interface (API) are unavailable offline. The recorded review could not use the candidate's browser route, so the candidate remains unpromoted and has no newer browser qualification. The exact four byte counts are retained in ``Evidence, scope and edition status''. They describe file packaging, with no measured interaction, loading-latency, memory, crash or GPU benefit.

The same separation applies to the observations on screen. Human-arm feedback is simulated. The human EEG/\allowbreak{}ESG/\allowbreak{}ENG/\allowbreak{}EMG preview comes from a separate evoked experiment on its own clock. Anatomical site, time, modality and missing pathways must remain visible so that it cannot appear to be simultaneous voluntary cortical-to-spinal-to-nerve control of the macaque-derived human-arm simulation.

\section{Alternative Explanations}

The distinction between a plausible display and a measured mechanism also changes how outside examples should be read. Useful approximations can succeed on some predictions and fail on others. Shiu and colleagues tested connectome-derived predictions through perturbations and behaviour. Specific successes coexist with missed Phantom and Usnea effects and incorrect G2N-1/Usnea water-silencing predictions. Missing basal activity and neuropeptide mechanisms---neuropeptides are small chains of amino acids used as signals---are possible limitations, not uniquely established causes. The evidence supports specified circuit predictions, not complete whole-fly fidelity.\citep{shiu2024}

The official March 2026 Eon account describes hand-chosen brain--body mappings,
heavy dependence on existing body controllers, missing state/plasticity and
unvalidated internal dynamics. It demonstrates an engineered integration under
its reported conditions, not independently recovered complete biological
function.\citep{eon2026} Pinned fly-chess and DoomFly sources add explicit
encoders, readouts or injected game-engine information; those software
capabilities cannot be attributed to living fly biology.\citep{flychess,doomfly}
These are source inspections and author disclosures, not reproduced gameplay
strength or proof of a lifetime-impossible biological ability.

Physics also needs independent checking. NeuroMechFly v2 reports slightly
below-ground simulated leg tips and a 0.05-millimetre margin used to prevent
false overstretch detection. It does not establish dramatic tunnelling or an
impossible FlyWire-brain capability.\citep{neuromechfly2024} Published first-hand
simulation accounts describe numerical or contact errors supplying artificial
energy to virtual creatures. A higher reward can likewise satisfy the wrong
objective, as in the reported CoastRunners case of circling reward targets
instead of finishing the race.\citep{surprises2018,specgaming2020}
These are different failure modes, not one invented universal error rate.

The same distinction matters when the output is a structure rather than movement. The original AlphaFold3 structural benchmark reports 62.9 percent success over 65 interface clusters from 71 complexes at DockQ greater than 0.23, with a 95-percent interval of 51.4--73.5 percent. DockQ measures reported structural-interface agreement.

Predictions were confidence-ranked from 1,000 seeds times five diffusion samples per target. Diffusion sampling generates a candidate structure by refining noise in stages; these are model samples, not biological replicates. The 37.1 percent complement belongs to that filtered benchmark and confidence-ranked sampling process. It is not a percentage of all proteins failing in real-world function. A structural interface score is not a functional assay.\citep{alphafold32024}

Crystallization limits some structure measurements, but cryo-electron microscopy and nuclear magnetic resonance provide other approaches. Crystallization alone supports no majority-unknown claim.\citep{pdbmethods} Pretraining on measured motion alone and declared unpaired transfer between homologous or mechanically analogous features with roughly similar biomechanics remain candidate training routes. Simulated or generated examples, including model-produced targets called pseudo-labels, may support training without each being a simultaneously measured neural--muscle--movement triple. Those examples remain generated or derived; agreement with the model or musculoskeletal simulation that produced them is not independently measured biological validation. Synthetic probes can check specified computational behaviour, but neither probes nor generated targets can supply a missing measurement. Biological claims still need comparison with independent measurements in the claimed domain.

The positive alternative is measurement. A brain--spine interface evaluated
cortical signals, spinal stimulation and actual standing/walking in one person
with chronic tetraplegia. Its measured interventional outcome is valuable
without becoming universal decoding or skill upload.\citep{lorach2023}
For this programme, predictive accuracy, mechanically plausible motion and
causal biological reproduction each require their own references and tests.


\section{Proposed Body and Teaching Tests}

A body-wide model needs a way to learn across bodies without pretending that unrelated macaque, human, fly and rodent recordings came from one nervous system. Shared functional roles---a limb, a joint, effort, contact and feedback---can support that proposed comparison, with separate verified body adapters. Homologous structures share evolutionary origin; mechanically analogous structures may serve comparable roles without that common origin. Either can motivate a transfer map when roughly similar biomechanics are specified. Anatomical similarity alone does not establish matching cells or interchangeable muscles.

Motion-only pretraining, unpaired representation learning and mechanical retargeting are candidate routes, alongside learning from paired neural and body observations. Scans can constrain geometry or test an anatomical assumption where they add information beyond the selected musculoskeletal model. For each transfer, state the source and target, joint coordinates, body frames, proportions, muscle roles, force and time scales, contact conditions, missing features and uncertainty. Measured, tracking-derived, model-estimated, retargeted, predicted and simulated quantities must keep those identities; transferred targets do not become simultaneous neural measurements.

Matched voluntary recordings, reliable clocks and body frames remain necessary for the claimed paired test across relevant levels, with missing measurements explicitly identified. Missing simultaneous observations limit that empirical claim; they do not forbid the proposed unpaired transfer routes. Compare proposed transfer with target-only training and naive or incorrect mapping under matched resources, preserving independent target evaluation and the original held-out restrictions. Virtual-rodent work offers a relevant example of combining an embodied model, movement and neural analysis, not a license to assume species transfer.
\citep{virtualrodent2024}

Both recurrent candidates fail the original no-harm and positive-neural-information requirements. All five added benefit claims fail too. Lower average error admits neither predictor. NeuroTranslation remains a proposed field; NeuroTranslator remains recording-specific software.

The proposed loop should expose each transition: supported observations become decoded estimates, pass through body adapters and return as feedback. Simple baselines must face the same matched tests. Illustrative motor replay and a deliberately goal-driven simulation can serve that limited question when their goals, assistance and simulator information are disclosed. A positive neural-control claim requires a predictor that passes its unchanged admission criteria. The current failed candidates do not pass.

Check the reward, the simulator and the independently measured task outcomes separately. Physical calibration, synchronized recording chains, participant-specific registration and generalization remain unverified. A controller that improves inside a simulation earns a claim about that tested computation. Biological teaching, retention and repair need their own evidence. Those missing biological results do not prevent a clearly limited computational proposal.

A later correction-learning comparison must keep the baseline and its information
resources fixed and test unfamiliar recordings where applicable, bodies, delays
and contact conditions. Simulator rewards for task completion, tracking, effort
and contact limits require verification against independently checked outcomes;
that verification is a future gate, not a result established here.

The biological teaching claim would require a further matched experiment with
unfamiliar tasks, appropriate controls and retained performance after trainer
instructions are withdrawn. Assistance budget and any necessary bridge must
remain controlled. Failure to beat a simple predictor, failure on a new
recording, or failure to retain performance must remain genuine ways for the
proposal to lose.

% R152 source proposal; provenance and exact prior text remain in the proposal manifest.
\section[A grounded conditional learner with distinct output interfaces]{A grounded conditional learner\\with distinct output interfaces}
Start with a movement we can measure well. For the proposed snap showcase,
the arm begins extended with the palm down. The elbow and forearm swing and
turn; the wrist and hand turn upward during the snap, then move down to rest.
The model must account for that whole movement, including the preparation,
contact and return. Finger pairing, organism-specific anatomy and timing still
need a real-reference series; a human-looking animation cannot establish that
a particular monkey can perform the same gesture.

The cup reach supplies a simpler precision comparison: relate recorded signals
to measured movement before attempting the coordinated snap. Each task still
needs its own anatomy, contact, timing and independent comparison.

That comparison connects three interacting groups: task-linked brain
observations, named-muscle observations and a musculoskeletal simulation.
A learner would relate the observations to the mechanics; independent movement
and contact measurements would test its predictions. EEG records scalp
electrical activity and surface EMG records muscle electrical activity.
A separately recorded nerve or spinal signal needs its own declared channel.
The proposed apparatus uses two arm sleeves, scalp EEG and synchronized,
calibrated left and right tripod views. Vendor, channel, bandwidth and clock
qualification remain open.

Intent to action is the target. By understanding and measuring the intermediary,
we aim to make that mapping as close to one-to-one as possible within a declared
task. Specify what counts as the intended action, the participant and session,
state and context, permitted history, tolerated error and independent reference.
The test then asks whether observing the middle reduces ambiguity and improves
prediction or qualified control. It does not assume a unique hidden intention
or claim that the best possible mapping has already been achieved.

These observations define what the learner would have to explain: recorded
brain and muscle activity, where the hand moved and when. Delivering sensory
feedback and producing an action would require further tests of separate
outputs; neither follows from a successful prediction. No paired trial, trained
model, delivered sensory output or physical actuation is established for this
proposal. Sources and permissions, physical and computational behaviour,
clinical interfaces, reference comparisons and held-out tests still need
independent assessment.

The retained LINK comparison answers a different question. Its 18 runs and
720 saved epochs are not the proposed paired EEG/EMG/movement experiment.
Both model-admission gates and all five benefit claims failed. Recording,
prediction, simulated movement, sensory delivery and physical action require
different evidence. EEG and EMG recorders alone cannot write sensation.
The prospective methods below state what the next experiment must measure.

\subsection{Form, function and action need a measured middle}
A structural map constrains what a system could do. A NeuroConnectome maps
specified neural connections; its organism, coverage, resolution and source
would have to be declared. The signals passing through those connections
still have to be observed or inferred, with the difference made explicit.
To predict the next action, we need the processes that turn structural
constraints into an outcome. The proposed programme follows form to function
to action, while intent to action names the task-level translation we want
to improve.

DNA to RNA and RNA to protein put a process between a starting constraint and
its downstream product. Regulatory context, RNA processing and cell state
change what that process produces; protein structure alone does not determine
the behaviour of a whole cell. AlphaGenome, RNA or cell-state models and
AlphaFold illustrate these different modelling jobs. AlphaTranscriptome names
a desired predictive capability here. Connecting the predictors would still
leave system dynamics and the environment to account for, with experimental
outcomes needed to test the combined prediction.
% Source-role pointers: alphagenome2026, scgpt, alphafold3; bounded owner checks.
\par\noindent \href{https://www.nature.com/articles/s41586-025-10014-0}{AlphaGenome regulatory prediction}~\cite{alphagenomeRetained};
\href{https://www.nature.com/articles/s41592-024-02201-0}{scGPT cell-state modelling}~\cite{scgptRetained};
\href{https://www.nature.com/articles/s41586-024-07487-w}{AlphaFold 3 structure prediction}~\cite{alphafold32024}.
These provide modelling examples, not evidence of an integrated cell programmer.
Our account of AlphaGenome and scGPT is limited to their reported modelling roles;
their complete methods have not been independently assessed here.

In the biological pathway, Brain2Spine and Spine2Body name proposed links to
test; their composition is the Brain2Body aim. The spinal cord integrates
descending and sensory signals and contributes local coordination, reflex
regulation and adaptation. This is a different mechanism from RNA processing,
but it serves the same explanatory purpose: account for the context-dependent
process between an initial constraint and its outcome. Supported spinal
observations would test that account. Where they are absent, the internal
states remain inferred or simulated. Physical spinal work remains read-only.
\par\noindent \href{https://pubmed.ncbi.nlm.nih.gov/38603479/}{Spinal sensorimotor adaptation}~\cite{lavaud2024}
is a constrained example, not evidence of unrestricted skill transfer or
validation of our proposed pathway. Our comparison draws on the abstract and
relevant author-manuscript passages; equivalence to the final published edition
and the complete methods and statistics remain unverified here.

Thought2Action relates a declared intention to an independently measured action.
Action2World follows the action to its consequence in the task environment;
Thought2World names the proposed composition. Body2World drew attention to the
consequence outside the body; Body2Action focuses on body mechanics and a
specified task outcome, while Brain2Motion focuses on the movement pathway.
These working labels identify related scopes for a test. They do not establish
interchangeable or working controllers.

The engineering route connects neural activity to musculoskeletal simulation
and then to a programmable model, toward NeuroTranslation. The proposed
movement-focused learner is a large movement model. System 3 thinking names
the proposed human--AI supervisory partnership around it. The supervisory
model would compare task strategies and diagnose where neural or mechanical
predictions miss independent measurements, then propose training or calibration
comparisons for the signal mapping and virtual body. The task learner would
perform the declared prediction or control; the supervisor would help choose
what to test. Neither role replaces the biological spinal intermediary or a
physical interface. Final held-out outcomes must remain outside rewards and
model selection.
\par\noindent Shaw and Nave's SSRN preprint,
\emph{Thinking---Fast, Slow, and Artificial}, and its
\href{https://executiveeducation.wharton.upenn.edu/thought-leadership/wharton-at-work/2026/05/thinking-fast-slow-and-artificially/}{institutional discussion}
provide conceptual context for this proposed human--AI role.~\cite{shawNaveSystem3Retained}
The comparison is limited to that framing; the complete preprint and institutional
article have not been reviewed here. Neither establishes a third anatomical brain
system or validates this programme's proposed learner or supervisor. Expert footage may supply
movement and tactical examples, with rights and camera uncertainty accounted
for; it supplies no measured expert EEG, EMG, muscle force or stimulation plan.

Action changes the outcome. Its sensory consequences inform perception and
the next intention. Independently measured movement would anchor the action
part of that loop; it would not by itself decode a unique sensation or
perception. Known stimuli, measured feedback and appropriate recordings or
reports would have to test those further mappings. The return path must
therefore be measured rather than drawn into existence. The full
Interoceptomics aim includes internal-body sensing and connectivity; external
touch and proprioceptive tasks contribute evidence but do not by themselves
solve that wider problem. A prediction, assisted execution and retained skill
without assistance answer different questions. Each wider claim must earn
its own independent comparison.

\subsection{Context connects. Grounding checks}
Within a reach, context has a concrete job: relate brain activity to muscle
activity in this task and at this time. Attention offers a candidate computation for those relationships, following
\href{https://arxiv.org/html/1706.03762v7}{Vaswani et al. (2017)}~\cite{vaswani2017} and \href{https://aclanthology.org/P19-1656/}{Tsai et al. (2019)}~\cite{tsai2019}. The model compares a query with keys and uses the resulting weights to combine their values. Masks keep unavailable observations and future
samples out. The calculation weights information; it does not establish that
attention understands all context.
% Source registry IDs: context-native-01a12f62, context-attention2017-v7, context-mult2019, context-policy4

A pattern can fit inside the model and still describe the wrong movement.
Grounding therefore means testing the declared prediction against qualified
observations under the same conditions. Independently measured movement would
check the outcome, while muscle activity would add an observation that
constrains the brain model. Even a matching reach cannot identify a unique
private experience. No acquired paired data or trained model is established
for this proposed loop; those comparisons still have to be made.

\subsection{Muscle activity tests the brain model}
Record brain and named-muscle activity during the same reach. In the proposed
paired trial, EEG would sample scalp electrical activity and surface EMG would
sample activity associated with the named muscles. Paired training could use
the muscle observation to constrain their relationship. Then withhold the
muscle target and movement reference from the EEG-only predictor: can the
permitted brain observations predict a specified muscle-activity observable
or movement?

Pairing and calibration remain proposed. Shared cues and movement artefacts
could give that predictor a shortcut, so the comparison needs controls for
both. The muscle observations would cover only the named muscles, and EMG
supplies neither afferent nerve activity nor direct force. Sensory claims
still need known stimuli and measured feedback. Several neural states could
remain compatible with even a well-matched pair; the pair could narrow a
task-specific explanation without revealing the complete hidden brain.
Checking the movement itself requires the next independent observation.
% Source registry IDs: context-native-01a12f62, context-policy4, overhaul-brantley2018
\par\noindent The proposed pairing still needs its own evidence. The candidate reference at \url{https://www.nature.com/articles/sdata2018133} has not yet been checked for exact support or reuse rights.

\subsection{Measured motion tests the virtual body}
Measure the same hand and arm from the aligned left and right tripod views.
The proposed motion-capture relay-composite would align the two point maps
in a common clock and coordinate frame, then join them into one composite
while retaining each point's origin. Before that joined map can check a
prediction, its calibration, alignment and capture precision would have to
be independently verified.

The movement map would constrain the external movement model: surface motion
and pose in the task's physical frame. To explain that motion mechanically,
the internal musculoskeletal model would describe specified bones, joints
and soft tissues. Organism, developmental stage and task would determine the
appropriate anatomical sources. X-ray projections could inform bone and
joint relationships; appropriate MRI references could inform specified soft
tissues. Registration, rights and source-specific limits would need independent
checks. No scan set is selected for this proposed comparison. Anatomy still
leaves activation, forces and neural function to test.

With a qualified composite, the musculoskeletal-simulation physics-engine's
joints, muscles, contacts and physical limits would face a measured comparison.
Fit the model on training data, then compare its predicted reach with
held-out measured movement. A force claim would require an additional force
measurement. The body's form limits possible motion; measured function would
test the mechanics and show where they need refinement.
% Source registry IDs: context-native-01a12f62, context-component-names-01a12f81, context-policy4, context-opensim2018, context-myosuite-docs, context-precision-01a12f77
\par\noindent \href{https://journals.plos.org/ploscompbiol/article?id=10.1371/journal.pcbi.1006223}{Seth et al. (2018)}~\cite{opensim2018}; \href{https://myosuite.readthedocs.io/en/latest/}{MyoSuite documentation}~\cite{myosuiteDocsUndated}. The simulation tools supply methods for that comparison; the proposed capture and body models still need its independent validation.

A lifelike hand makes that comparison visible. It cannot establish biological
fidelity. Capture alignment, precision, mechanics and biological fidelity have
not been demonstrated by this proposal, and captured geometry cannot establish
hidden neural state, internal forces or unique subjective experience. The
measured reach would check a movement prediction within those limits.

\subsection{Two pairs. One conditional learner}
The two comparisons would meet in the proposed large movement model: EEG
observations paired with muscle activity, and measured movement compared
with predicted mechanics. Large movement model names the proposed family;
deep neural network describes a possible implementation. Large refers to
the intended scope, rather than a demonstrated parameter count. Encoders
would receive only the observations permitted for the named test. Two
branches could use the attention and multimodal computations of
\href{https://arxiv.org/html/1706.03762v7}{Vaswani et al. (2017)}~\cite{vaswani2017} and \href{https://aclanthology.org/P19-1656/}{Tsai et al. (2019)}~\cite{tsai2019} to summarize those relationships
in a shared task representation; separate output parts would predict each
direction's targets.
% Source registry IDs: context-native-01a12f62, context-policy4, context-attention2017-v7, context-mult2019


Signal features belong to branches, and branches belong to the learner.
That hierarchy describes the computation, not two anatomical thinking systems.
A common clock, task and physical frame would connect the branches; independent
measurements would check their errors. Form constrains the prediction, and
checked function would guide refinement. Several hidden states could still
fit the same reach, so neither a shared representation nor conditional
predictions would establish an exact inverse nervous system. The architecture
has not been selected, implemented, trained or independently accepted.
To choose it, we would have to compare it with simpler baselines on
appropriately reserved observations.

\subsection{Learn the movement, then test what transfers}

A reach can provide a training target even when no EEG accompanies it. Earlier movement supplies the input; a later observation supplies the target. A separate proposed comparison would ask whether learning that relationship in one person helps forecast another person's grasp-and-lift movement. The WAY-EEG-GAL descriptor provides a relevant continuous-recording candidate, alongside its distinct windowed trials and derived lift summaries.\citep{luciw2014way} This movement-only question needs no AllLifts Run/trial join or invented cross-person neural pairing. The original condition-specific five-weight study remains a separate design.

The mapping has to explain what two bodies share. Homologous structures share evolutionary anatomical origin; mechanically analogous structures can serve comparable roles without that origin. A human digit aperture and a robot gripper opening could support a task-space comparison after separate mechanical qualification; this human-only proposal would not test that robotic extension. Joint degrees of freedom, axes and signs, proportions, muscle--tendon roles, loads, inertia, contact and time scales remain body-specific. Similar geometry cannot establish equivalent dynamics or neural coding. Retargeted motion remains derived; tracking, estimation, prediction and simulation keep their different evidential roles.

For the proposed two-person comparison, the representation contains index-, thumb- and object-sensor positions relative to the wrist, digit aperture and causal wrist velocity: thirteen features, with thirteen availability masks. Each participant needs qualified position units, a proper wrist-frame rotation, orientation convention and sensor-attachment calibration. A positive median wrist-to-index-sensor distance, expressed in qualified metres, supplies the characteristic length from assigned training support only; it is not a bone measurement. Multiplying normalized velocity by the same fixed 0.1-second reference interval makes all thirteen features dimensionless. Physical velocity errors invert that scaling with length divided by 0.1 seconds. Missing features receive no target loss; an unavailable observation is not a measured zero. Full finger joints, individual muscle forces and subject-specific inertia are not supplied by these sensor features.

The continuous target would be the same observed feature set 100 milliseconds beyond the latest actually available filtered input time. The proposal retains fifty frames at 100 Hz and a causal 51-tap filter before decimating qualified 500-Hz kinematics by five. Its 50-millisecond linear-phase delay must remain in the physical timing account; shifting information backward would change the prediction claim. Coefficients, passband, timestamp convention, gaps and boundary handling still require qualification. Only named kinematic columns may enter: the documented kinematics matrix also contains force-plate signals.

Identity and independence answer different questions. Trustworthy versioned participant-release attribution, an immutable selected archive entry and a qualified capture can identify the participant record and input object without requiring one particular internal ID-field spelling. They cannot establish biological identity or independent original acquisitions. Shared samples, trims, derivatives and overlapping recordings must form indivisible acquisition groups before filtering, windows, normalization or role assignment. Different names, checksums or series numbers do not prove separation. The complete public non-ST candidate roster and these relations still need establishing; all 294 rows of the prior AllLifts table remain excluded from training, evaluation and assigning trials or acquisition groups, and ST, competition storage and the reserved LINK recordings remain unavailable.

Qualified P1 groups would supply source pretraining; P2 would supply target training and development. With $G\geq2$ independently established P2 acquisition groups, the fixed identifier-hash ranking assigns the first $\lfloor3G/5\rfloor$ to training and the remainder to development. Canonical original-parent identifiers, their serialization and the roster must be fixed before ranking or examining outcomes. This is a nominal 60/40 split, not chronological ordering or a promise of condition balance. If unresolved ancestry merges P2 into one group, the comparison stops; splitting its windows cannot manufacture an independent development group. No group or segment has been assigned by the proposed identity method.

Three matched arms would test the transfer. B0 pretrains on correctly mapped P2 training movement. B1 and B2 use identical P1 windows and order: B1 sets the index-role input to the thumb sensor, the thumb-role input to the object sensor and the object-role input to the index sensor, while B2 retains the correct correspondence. Both share the correct source calibration and scale before that corruption, keep the wrist anchor, and use the same correct target adapter and ordered target batches. Parent-balanced sampling must record unique windows and repetitions. The false mapping can be relearned or can alter feature availability; matched update counts do not match source diversity. These are limits on attributing any difference to anatomical correspondence, not reasons to replace the control after seeing its result.

The proposed learner is one GRU layer of hidden width 32, with 26 inputs and thirteen outputs. Each arm and fixed seed---17, 31 and 73---receives 200 pretraining and 100 target updates, batch size 16, identical per-seed initialization and the fixed last checkpoint. Adam uses learning rate 0.001, betas 0.9 and 0.999, epsilon $10^{-8}$ and zero weight decay, without clipping, scheduling or early stopping. Training minimizes mean squared residuals over observed target features; a batch with no observed targets is rejected. The resulting 2,700 updates are planned, not performed. No external checkpoint is selected; numerical updates, filtering, precision and actual workload limits remain to be checked independently.

For each development parent and seed, the proposed score takes the square root of summed squared errors divided by the count of observed target features, then averages parent scores without weighting by parent duration. A zero denominator rejects the comparison. Report feature counts, missingness, physical-unit errors, uncertainty and persistence at the same information endpoint, including every worse parent or feature. B2 would have to lower normalized error against both B0 and B1 in every fixed seed, with no applicable original criterion failing, to justify another investigation. This remains a two-person development comparison, not independent final confirmation, population transfer or biological teaching.

Broader motion collections could supply additional movement examples only after their exact assets, rights, task coverage and original-recording overlap are qualified. In \href{https://www.addbiomechanics.org/download_data.html}{AddBiomechanics}, measured markers and ground reactions must remain distinct from fitted kinematics and estimated inverse-dynamics torques. \href{https://amass.is.tue.mpg.de/license.html}{AMASS} requires checking the selected asset's research-use, non-commercial and redistribution restrictions. Copies across collections stay with their original acquisition. LaBraM's WAY source-family exposure is flagged, while exact checkpoint/participant/trial exposure remains unknown; reserving a recording later cannot erase earlier pretraining exposure. No such checkpoint is selected for this comparison.

Before adding X-ray/CT or MRI, inspect the anatomical sources already used by the selected musculoskeletal model. A scan is useful when it constrains a stated gap or supplies an independent check beyond that model. Its person or species, scale, registration, rights and uncertainty still matter. Another person's scan remains an external anatomical reference; a generic model does not implicitly measure this participant's geometry, activation or forces. No scan set is selected here.

Actual source eligibility, acquisition groups, units, frames, clock and calibration remain unresolved for this continuous comparison, so no new fit or successful transfer is reported. A useful movement prior could later be tested with actual paired neural/body observations and the permitted input masks. Its transferred targets would still not become simultaneous measurements. The next result must identify where the mechanical forecast improves, where it worsens and which measurement or assumption accounts for the difference; neural equivalence, causal physiology, sensory delivery and retained skill would require further experiments.

\subsection{Pretrain. Pair. Verify. Refine}
Before a paired example can teach the mapping, we have to know whether we may use it and what it measures. Data rights, calibration and the held-out subjects, sessions and tasks determine which observations may enter; resolve those conditions before processing data or choosing a model. Open-source code supplies no blanket permission for datasets, checkpoints or clinical interfaces. Qualified data could first train a representation for each recording type, then synchronized paired observations could train the declared mapping from permitted inputs. LaBraM is one candidate EEG encoder; it does not supply a complete brain-body writer.

Paired prediction would give later refinement a baseline to beat. Before
optional reinforcement learning, post-training or task-specific fine-tuning,
the physical predictions and reward need checking against independent
outcomes. Offline reinforcement learning needs qualified coverage of actions,
transitions and rewards: labels alone cannot tell us what an untried action
would do. Backpropagation can update differentiable neural components, but
it cannot automatically pass through every physics engine; universal
differentiability is not established. Simpler baselines and frozen final
tests remain through every phase.
% Source registry IDs: context-native-01a12f62, context-policy4, context-labram2405, context-cql2020-abstract, context-opensim2018, context-myosuite-docs
\par\noindent \href{https://arxiv.org/html/2405.18765v1}{Jiang, Zhao and Lu (2024)}~\cite{labram2024}; \href{https://arxiv.org/abs/2006.04779v3}{Kumar et al. (2020)}~\cite{cql2020}; \href{https://journals.plos.org/ploscompbiol/article?id=10.1371/journal.pcbi.1006223}{Seth et al. (2018)}~\cite{opensim2018}; \href{https://myosuite.readthedocs.io/en/latest/}{MyoSuite documentation}~\cite{myosuiteDocsUndated}. These sources offer candidate methods for the sequence; validation of the proposed paired system still requires the experiment below.

No data acquisition, training, fine-tuning, RL rollout or accepted checkpoint
is established for this proposed loop. The next experiment must test what
the prediction and reward represent against independent outcomes. A reward
that the computer can calculate supplies no biological proof on its own.

\subsection{Respond quickly. Improve carefully}
A response should not become a new experiment every time it is issued.
The proposed response path would use a frozen accepted model to predict from
permitted past and present observations and prepare encoded outputs for a
qualified interface. A separate path would inspect errors, calibrate and
propose a candidate update. Independent evaluation and a restorable checkpoint
would precede promotion. Updating an estimate of the current state would
therefore remain distinct from changing learned weights; examining errors
would not silently change the running model.

This fast-and-slow comparison is a Kahneman-like analogy for the proposed division of work. It does not identify two biological systems or assign inherent thinking speeds to EEG and spatial data. Real-time performance would have to be measured from acquisition through filtering, alignment, inference, output, feedback and rendering, including delay tails and jitter. No accepted model, real-time demonstration or running system that adapts under these rules is established here. The loop has no measured end-to-end latency or accuracy guarantee.
% Source registry IDs: context-native-01a12f62, context-policy4

\subsection{Write sensation. Enact action}
The prediction ends at the computer. A separate interface would have to turn
it into sensation or action. Input and output describe direction relative to
a computer; afferent and efferent describe physiological direction relative
to the central nervous system. A computer output could feed a sensory
transducer, or a controller and actuator for a virtual, robotic or later
qualified biological body.

A display or haptic interface supplies sensory feedback, which differs from direct neural writing. Direct cortical writing could have a sensory role without being a peripheral afferent signal, just as a robot command is not an anatomical efferent nerve. EEG and EMG recorders do not write, and predicting an EEG-like pattern does not deliver sensation. Each proposed output needs an independently qualified interface, permission for its use and a measurement of what it delivers. No such physical interface, stimulation, acquisition or actuation is established by this proposal.

Comparing a recorder with a sensory or action device requires evidence about what each measures or delivers. The clinical source could not be read, so it supplies no methods evidence here. Implanted precedents cannot establish that EEG or EMG recorders are equivalent to devices that deliver sensation or produce an action. Those physical additions remain unpriced and outside the acquisition and virtual-modelling scope.
% Source registry IDs: context-write-native-01a12f66, context-native-01a12f62, context-policy4

\subsection{Close the loop with evidence}
Return to the cup and test each prediction against the observation it was
asked to explain. The brain-to-body prediction would face withheld observations;
the body-to-recorded-neural estimate would face its own recorded target.
A sensory or action interface would need a further test of what it delivered,
under which conditions, with what error and delay. Two models can agree and
still be wrong; a high reward can have the same problem. Neither agreement nor
a successful prediction would, by itself, identify a unique inverse brain or
establish someone's private experience.

The programme's distinction remains that approximation gets us towards plausibility and virtualization gets us to precision. Independent measurements, repeatability, causal checks and stated error bounds must earn that precision. It is not an automatic result, and the proposal establishes no perfect biology. Keep the original held-out, real-reference and retained-learning controls, including the distinction between useful guidance and learning that remains after guidance ends.

Acquisition and virtual-modelling costs do not establish the cost of delivering
sensation or action. Physical sensory writers, transducers and actuators remain
unpriced additions. No loop, training or
physical-writing result is claimed here. The next experiment must report where
each declared prediction works and where it fails against its independent
measured target.
% Source registry IDs: context-native-01a12f62, context-write-native-01a12f66, context-precision-01a12f77, context-policy4, context-spec15, context-glossary13, b2b-circular-checklist-original

\subsection{Technical representation and mask}
Learned query, key and value representations give the proposed attention
calculation a precise form. An admitted mask would restrict which observations
may contribute; biological grounding still requires the independent comparison.
\[\operatorname{Attention}(Q,K,V)=\operatorname{softmax}\!\left(QK^{\mathsf{T}}/\sqrt{d_k}+M\right)V.\]
Here Q, K and V are queries, keys and values; M expresses the admitted mask. The dimension $d_k$ is the common feature width of each query and key vector. The equation is methodological context from the supplied technical proposal. Modality, channel, task and time are part of the observation contract. No weights, EEG encoder or gradient route are selected or executed.

\subsection{Permitted observations and causal modality routes}
Declare B=EEG, E=EMG, Y=independently measured movement, C=task/\allowbreak{}stimulus/\allowbreak{}context, and P=the source-bound physics model. Training may relate qualified paired observations. Inference uses only its explicitly admitted subset:

\begin{itemize}
\item EEG-only: permitted past/present B and declared context predict a specified   EMG/\allowbreak{}movement target. Withheld E, Y and future samples cannot enter the predictor.
\item Body-only: permitted E and/or Y predict specified recorded neural features as   a conditional distribution with uncertainty. This is not a unique inverse brain.
\item Hybrid: additional modalities require a separately named route and evaluation;   a hybrid result cannot be reported as EEG-only.
\item Simulation-derived inputs stay simulated. They do not replace independent   observations or inherit the biological validity of a rendered body.
\end{itemize}

Separate training signals from evaluation references. Cycle consistency can make two models agree while both are wrong. Shared cues, movement artefacts, subject identity, session overlap and preprocessing leakage need the original controls. Do not expose final held-out targets to model, physics or reward selection.

\subsection{Representation, supervised learning and policy refinement}
\begin{enumerate}
\setcounter{enumi}{-1}
\item Establish source rights, subjects/tasks, measurement units, calibration,    uncertainty, clock alignment and frozen subject/session/task splits. Audit    pretraining overlap before preprocessing and model selection.
\item Pretrain modality representations on the qualified inventory. LaBraM is one    candidate EEG encoder: neural-tokenizer training, masked neural-code prediction    and downstream fine-tuning. It is not a complete brain-body writer. Code,    checkpoint, dataset and individual-use rights remain separate qualifications.
\item Learn paired supervised mappings with named targets and explicit modality    masks. Compare appropriate simpler, cue-only and artefact baselines. Keep raw,    processed and model-derived values distinct.
\item Qualify physical predictions and the reward verifier against independent    measurements. A successful task, smaller calibrated movement error and valid    physical states are prospective reward ingredients, not completed results.    Declare loss/reward weights on training/development evidence. A computable    objective or high reward does not confer biological ground truth.
\item Where admitted action/transition/reward coverage supports it, evaluate optional    RL and post-training. Offline trajectory labels alone do not establish the    consequences of actions never supported by the data. Conservative Q-Learning    is a methodological candidate, not a safety or performance guarantee here.
\item Fine-tune for the declared task and qualify both directions on frozen final    tests, uncertainty, baselines and every original threshold. Keep failures.
\item Serve a frozen accepted version. Propose later updates separately and promote    only after independent acceptance, restoration proof and applicable authority.
\end{enumerate}

An illustrative supervised objective would combine declared errors for
specified EMG, movement and recorded-neural targets with applicable physical
constraints and optional consistency terms. Each term needs its units,
measurement uncertainty and target definition before a weight is chosen.
Consistency can check whether the outputs fit together; an independent
observation checks whether they fit the measured outcome.

Backpropagation changes differentiable components of the computer model;
biological feedback is a different process. A physics engine in the path does
not automatically give the learner a gradient. Qualify the actual
OpenSim/MyoSuite rollout and optimization route, or a separately validated
surrogate, rather than assuming gradients through every route. No training,
RL rollout or task checkpoint is created or certified here.

\subsection{Frozen response and independently accepted adaptation}

What must the response path preserve when an observation arrives late? The frozen-weight distinction in ``Respond quickly. Improve carefully'' means that the state estimate may change as observations arrive, while learned weights change only after a separately tested candidate is accepted. Each response needs its causal time window, permitted observations, target and output preparation. Future observations cannot repair a missed deadline, and an unaccepted candidate update cannot become the running model.

The timing test must follow that response through the whole chain specified there, including transduction or actuation where applicable. Delay tails, jitter, missed deadlines, uncertainty and coexistence with healthy workloads determine whether the chain can meet its declared task; component inference time cannot settle that question. The Kahneman-like analogy describes this division of computational work. It establishes neither two anatomical systems, inherently fast or slow sensor modalities nor real-time performance.

\subsection{Distinct sensory and action interfaces}

The output tests follow the two routes in ``Write sensation. Enact action''. A visual, haptic or other sensory transducer would be tested by stating the sensory content, interface and authority, then measuring the actual response. Direct neural writing remains a separate later research interface; a sensory role would not automatically make it a peripheral afferent signal. The boundary matters: a recording related to the body's motor output can enter a computer as input. Computer direction and direction relative to the central nervous system must therefore both be stated.

For the action route, the test would identify the control delivered to the virtual body, robotic actuator or later qualified biological interface and measure the resulting action. The robot command remains a computer output, not an anatomical efferent nerve. That distinction determines the evidence needed for simulated action, physical actuation, retained learning and biological causation. Predicting an EEG-like pattern delivers no sensation, and EEG and EMG recorders supply neither of these writing functions.

The output design therefore needs explicit subinterfaces within the unchanged five programme pillars: sensory encoders and transducers for sensation, and action encoders, controllers and actuators for enaction. This document gives no stimulation dose, implantation, electrode-placement, acquisition or actuation procedure, and grants no participant, clinical, device, purchase or deployment permission.


% Proposed insertion only. No master/include/source selection is changed.
% Inherit the selected master's Noto Sans/Fira Math setup; no new font defaults.
\section{More data must survive a harder test}

A thousand repetitions of one reach still describe one kind of reach. A movement model trained across people, recording days, tasks and bodies could encounter relationships that a narrower collection never contains. Whether that breadth improves a particular prediction is the scaling hypothesis proposed here. Increasing the amount of qualified data and increasing its diversity are separate interventions. Neither guarantees improvement, and the proposed Large Movement Model has not been trained or validated as a body-wide world model.

\subsection{Four source groups, distinct observations}

The proposed training programme draws on at least four groups: neural observations; muscle and other myographic observations; measured motion and displacement; and data generated by physics-based motion or musculoskeletal simulation. These are source groups, rather than a claim of exhaustive coverage or four fully acquired corpora. The neural scope includes suitable brain, spinal, peripheral and relevant visceral observations, with their recording sites, modalities and task roles stated separately. The present emphasis is motor--sensory. Myoception is the programme's term for the sensory side of musculoskeletal contraction and movement, including proprioception, tactile feedback and other task-relevant sensations linked to action. It marks proposed programme scope, without establishing a new external scientific standard or a body-wide simultaneous recording.

Tactile stimulation, proprioceptive feedback, contact, force and task context require different observation relationships. A known tactile stimulus records what was delivered; a contact sensor records its specified contact event; a force instrument measures its calibrated physical quantity. Kinematics may help estimate body configuration without directly recording sensory afferents or the subject's proprioceptive experience. EMG supplies muscle electrical activity, not direct peripheral nerve or sensory recording. Measured stimuli, reported sensation, recorded neural features and model estimates must keep their identities even when they describe one action.

The fourth group contributes training examples as well as physical constraints. A qualified generator could supply simulated trajectories, state transitions, muscle--tendon or contact-force outputs and declared movement priors under specified controls. Every example needs its simulator and body-model version, parameters, solver, control representation and source lineage. A joint-target controller and a muscle-driven controller expose different action spaces; one must not be reported as the other. Motion retargeted from a recording retains that recording's lineage and the transformation, while simulated states remain simulated. A generated force is not an independent force measurement.

Before using such examples, calibrate the generator's relevant parameters against suitable measured support, retaining calibration uncertainty and training-only role boundaries. Compare a fixed-real-data reference with a simulation-augmented arm under matched permitted information, compute and original criteria. Record real-example exposure, simulated exposure and updates in each arm: a fixed real corpus alone does not match training exposure. Report simulated amount, simulated diversity and real acquisition amount separately; then test the synthetic contribution against held-out real observations. Keep those real tests out of generator tuning, calibration, reward choice and model selection. Errors in contact, dynamics or source correspondence must remain ways for the synthetic contribution to fail. Generated data cannot supply independent biological ground truth for its own assumptions. The generator's methods, implementation and use conditions still need checking; no frontier implementation or dataset is selected here.

The four source groups and five component roles describe different axes. Brain signals, muscle signals, measured movement, a virtual body and a learning model describe jobs in the proposed system. A virtual body can both generate training examples and constrain forecasts; the learning model can use several source groups. Neither arrangement closes the inventory to other useful, properly measured information.

\subsection{Separate amount from diversity}

For each modality, training amount must identify what was counted: original acquisition groups, people or animals, recording duration, trials, observed target values and distinct training windows. Overlapping windows and repeated passes do not create independent acquisitions. The study must also report diversity as a description of the training distribution: organism, subject, session, task, body region, recording modality, apparatus, environment and measurement quality. One additional subject and one additional sensor are different changes. A count of datasets cannot make them interchangeable.

The proposed amount comparison would retain a declared source distribution while increasing eligible training support. Its diversity comparison would retain a declared amount while changing specified dimensions of that distribution. A crossed design could estimate their interaction. All comparisons would keep the architecture, allowed input information, target, preprocessing and model-selection rules fixed. Training exposure, repetitions, updates and measured computational cost must be reported together: equal update counts do not by themselves establish equal compute or equal information. Where different modalities or recording quality prevent a matched comparison, that limitation becomes part of the result rather than a hidden explanation for it.

Each new comparison needs its own prospectively frozen roster, budgets, hypotheses, seeds, controls and decision rules. It does not replace the adopted continuous WAY-EEG-GAL development design or alter its fixed source/target roles, updates or selection. Nor may it reopen the reserved LINK dates, ST or other protected tests. Training-only calibration and normalization follow original acquisition groups through every trim, derivative, alias and window. An exposed recording cannot become a fresh test by changing its name.

Generalization must name the change being tested. An excluded session tests a different question from an excluded subject, task, apparatus, dataset or species. The proposed evaluation would report these domains separately, with a fixed target-only/no-transfer reference, the appropriate incorrect-correspondence control, simpler predictors and the original null, shuffle, confound and causal-time controls. New data may help familiar cases while harming an unfamiliar body or task; those failures and any negative transfer must remain visible. A pooled improvement cannot overrule an original no-harm requirement.

\subsection{A power law is a candidate description}

For one defined error measure, one evaluation population and one fixed diversity condition, a candidate relationship is
\[
E(N)=E_{\infty}+A\left(N/N_0\right)^{-\alpha}.
\]
Here $N>0$ is the registered measure of training amount, $N_0>0$ is a fixed reference amount in the same units, and $E(N)$ is a specified nonnegative held-out error. For this decreasing-error candidate, require $E_{\infty}\geq0$, $A\geq0$ and $\alpha>0$. $E_{\infty}$ and $A$ have the error's units; the ratio and exponent are dimensionless. With $A>0$ the curve decreases strictly as $N$ increases; $A=0$ gives constant error. A positive exponent alone is insufficient. Any upper bound of the chosen metric must also be respected throughout the declared fitting and prediction range. The fitted floor would be a parameter of this candidate, not an independently established irreducible biological error. This form would establish neither a universal law, unlimited extrapolation nor improvement from every additional dataset.

The form is a hypothesis, with no fitted exponent reported here. Its decreasing family does not constrain what the observations are allowed to show. Constant performance, saturation, increasing error, non-monotonic change and domain-specific deterioration must be credible alternatives, and negative transfer must be reported even when this curve cannot describe it. The fitting range, alternative models, uncertainty procedure and predictive check at withheld scaling levels must be fixed before inspecting the relevant outcomes. Related windows cannot serve as independent uncertainty replicates. An apparent curve driven by one additional subject, a changed target definition or a less demanding test distribution must be diagnosed at that level.

The errors themselves retain their meanings. A normalized EMG-target error, an image-pixel displacement error and a physical joint-angle error cannot be pooled into a biological accuracy percentage without an independently justified scoring rule. Report the original units, per-domain errors, sample denominators, missingness and uncertainty alongside any declared dimensionless summary. Acquisition diversity, computational model size and statistical averaging remain separate axes.

\subsection{Complementary recordings leave missing links}

FALCON M1-A offers a concrete neural--muscle candidate: the fixed DANDI release describes simultaneous M1 activity and preprocessed intramuscular EMG from one rhesus macaque. Its three split labels do not represent three animals. No released synchronized motion-capture trajectory is established by the reviewed metadata. The published target processing is acausal, so an offline result using that target would not establish a purely causal real-time system. Its small minival subset comes from the calibration data and is used to check submission format; it is not an independent validation population. These constraints belong beside any proposed use of the dataset. \href{https://doi.org/10.48324/dandi.000941/0.241029.1405}{FALCON M1-A release}; \href{https://papers.nips.cc/paper_files/paper/2024/file/8c2e6bb15be1894b8fb4e0f9bcad1739-Paper-Datasets_and_Benchmarks_Track.pdf}{FALCON methods}.

AJILE12 offers a different relationship: intracranial recordings and video-estimated, two-dimensional upper-body pose across twelve clinical participants. Its pose coordinates are image pixels, not instrumented three-dimensional motion capture; paired EMG is not established in the reviewed release. Pose-label construction is participant-adapted, and possible exposure of a particular future held-out day to the pose estimator remains unverified. The numerical alignment uncertainty also remains to be qualified. A neural predictor's unseen subject or day therefore does not automatically make its target-label construction independent. \href{https://doi.org/10.1038/s41597-022-01280-y}{Peterson et al., 2022}; \href{https://doi.org/10.48324/dandi.000055/0.220127.0436}{AJILE12 release}.

These sources provide complementary candidate links. A macaque's neural--EMG example and a person's neural--pose example never become one simultaneous neural--muscle--movement observation. Qualified unpaired pretraining could learn source-specific representations; an explicit transfer model could then be tested against actual target measurements. It would still have to earn its cross-domain correspondence. These release descriptions establish possible source relationships, rather than downloaded recordings or a trained mapping. The exact use conditions, file integrity, clocks, units, earlier exposure and independence of the chosen partitions still need checking before training.

\section{A world model must account for the body it moves}

Consider a hand closing on an object. The exterior changes: fingers displace, contact begins and the object may move. Inside the model, muscle--tendon forces and joint motion must account for that change under the declared load. The proposed world model joins those questions. Its interior initially means the specified musculoskeletal states and mechanics; its exterior means the supported body motion, objects, contacts and environment. The phrase does not imply that every organ, neural pathway or private sensation has been reconstructed.

NeuroTranslate would supply conditional mappings between declared observations and targets. MoCap ReComp names the proposed source-preserving composite movement reference; a video-derived pose collection does not thereby become that instrument. MyoSim Replay names the proposed physical simulation component. Their proposed combination is the Large Movement Model. A neural composite here would retain identified recording types and source-conditioned representations. Combining them in software would not identify equivalent neurons across people or species, or assemble a measured whole nervous system.

\subsection{Keep observations and internal state distinct}

A prospective model could represent mechanical state by joint configuration $q$, generalized velocity $v$, muscle-model activation $a$ and declared environmental state $c$. For a rotational joint, $q$ and $v$ may use radians and radians per second; translations require their own length units. The simulator's state dimensions and conventions determine the actual representation. A learned transition would predict a distribution over later states, conditioned on the permitted past observations, task context and any declared control input. At inference, withheld measurements and future samples remain unavailable even if they supplied training targets.

Each sensor needs a separate observation relationship. A neural recording measures its named electrical or imaging quantity. EMG measures muscle electrical activity; its processed amplitude is not automatically calibrated muscle activation or force. Tracking supplies observed or estimated geometry with its calibration and uncertainty. A simulated muscle force remains a model output unless an independent force reference tests it. Missing observations receive an availability mask, not a fabricated zero or a synthetic measurement label. The reverse prediction of recorded neural features remains conditional and potentially non-unique.

For a declared articulated-body model, one proposed mechanical check can be written schematically as
\[
M(q)\dot v+h(q,v)
 =R(q)F_{\mathrm{mt}}+J_c(q)^{\mathsf T}f_c+\tau_{\mathrm{ext}}.
\]
$M$ is the model's inertia operator and $h$ collects its declared velocity and gravitational terms. Muscle--tendon forces $F_{\mathrm{mt}}$ act through the moment-arm or generalized-force mapping $R$; contact forces $f_c$ act through contact Jacobian $J_c$; $\tau_{\mathrm{ext}}$ contains other declared generalized forces. Each equation component must use the model's compatible coordinates and generalized-force units, with the coordinate-to-velocity mapping stated where $v\neq\dot q$. This is a proposed consistency relation for a specified simulator, not an inferred equation of the participant's complete biology. Geometry, inertia, muscle parameters, external loads and contact assumptions still require independent support. OpenSim provides established musculoskeletal modelling methods, rather than validation of this proposed assembly. \href{https://journals.plos.org/ploscompbiol/article?id=10.1371/journal.pcbi.1006223}{Seth et al., 2018}.

Predicted movement must be compared in aligned time and coordinates with independent movement references. Internal force claims need appropriate force measurements; matching the visible path alone cannot identify the force distribution that produced it. Test contacts, joint limits, energy accounting, gravity, friction and solver sensitivity within the claimed domain. A model that exploits a numerical correction or receives privileged future targets can improve its score without improving the intended forecast. Report that failure as such. A more lifelike rendering cannot repair it.

The first useful result may be a limited forecast with explicit missing state, rather than a complete body. Its tests would isolate the neural representation, muscle observations, movement prior and physics constraint through matched ablations and external outcomes. Any later policy refinement would need suitable observations of actions, their consequences and rewards, with its reward calculation checked against those independent outcomes. Predictive fitting, offline replay, simulated control, physical intervention and retained biological learning remain different experiments.

% R4 insertion candidate only; no existing manuscript is modified.
% Inherit Noto Sans/Noto Sans Mono/Fira Math from the selected master.
% Requires the shared result table and the proposed supplementary bibliography.
\subsection{One muscle pulse, two simulated trajectories}

A muscle command should change the body the simulator computes. We tested that limited relationship in one installed human right-hand mechanism, \texttt{myohand\_r}, attached to a fixed passive torso. The question was whether a specified excitation changes the recorded mechanical state under a matched comparison. No neural recording, measured EMG or measured movement supplied the input or reference. This is a numerical software diagnostic within the wider mechanics proposal; it does not validate a biological hand or select a training generator. The model came from \texttt{myo\_sim} 0.2.3 at commit \texttt{93b0ca8f4ec90c9899ee7f05fee561e9911da91b}, with MuJoCo 3.14.0 and NumPy 2.5.3.\cite{myohandMyoSimPinned2026,myohandMuJoCo3140,mujoco2012,myohandNumericalRecord20261105}

\subsubsection{Mechanism and matched comparison}

The compiled mechanism has 23 scalar hinge degrees of freedom, 39 muscle--tendon actuators, 39 activation states and zero sensors. Its coordinates are joint angles in radians. Both branches began from matching full recorded initial states, including joint positions, velocities and activation states. The control branch received no additional excitation pulse; the other received dimensionless excitation 0.2 at the actuator named \texttt{FDS2\_r}, with all other controls zero. This input is a muscle-model command, distinct from a joint-position target or an electrical recording.

Both branches used the native Euler setting with $\Delta t=0.002\,\mathrm{s}$ for 100 explicit steps, giving a nominal horizon $T=0.200\,\mathrm{s}$ per branch. Pulse steps 10 through 59, counted from zero, use the pre-step interval $0.020\leq t<0.120\,\mathrm{s}$. Gravity and passive dynamics remain active in the comparison: zero added excitation need not leave the hand still. Initial and final states were each synchronized with an explicit forward evaluation. The record therefore contains 200 explicit integration steps and four explicit forward evaluations. Two compilation calls and one internal construction-time forward evaluation are inferred from the executed source path; these are not counts of every internal native solver operation.

Construction also matters. Although the optional \texttt{inertia\_floor} argument remained \texttt{None}, the native builder changed each diagonal inertia component of four fixed scaffold bodies from $1$ to $10^{-4}\,\mathrm{kg\,m^2}$. That transformation belongs to the tested mechanism and must remain with its provenance. No additional physical-parameter retuning or fallback to \texttt{implicitfast} was used. Position, orientation, model scale, contact behaviour and biological fidelity have not been independently calibrated by this diagnostic.

\subsubsection{Recorded response}

Let $P$ denote the pulse branch and $C$ the matched control. At the synchronized final state, define
\[
\Delta q_{\infty}=\max_{1\leq j\leq23}|q_{P,j}(T)-q_{C,j}(T)|,
\qquad
\Delta a=a_{P,\mathrm{FDS2}}(T)-a_{C,\mathrm{FDS2}}(T),
\]
\[
\Delta F=F_{P,\mathrm{FDS2}}(T)-F_{C,\mathrm{FDS2}}(T).
\]
Here $a$ is the selected model activation, with no physical unit, and $F$ is the selected signed tendon-actuator force in newtons under the model's convention. The coordinate maximum is across the 23 hinges at the final time, rather than across the whole trajectory. The force uses the final synchronized forward evaluation; a force retained from the last integration step need not refer to that same state.

% Shared by all four proposed R3 inserts; no added packages or font defaults.
\begin{table}[htbp]
\centering
\caption{Final contrasts for one simulated right-hand mechanism. All values compare the pulse branch with its matched control at nominal $T=0.200\,\mathrm{s}$. The angle value is the largest absolute separation across 23 hinges at that endpoint; activation and force retain the pulse-minus-control sign. Activation and force concern the selected \texttt{FDS2\_r} actuator. Force uses the synchronized final forward evaluations and the model's sign convention. These are simulated contrasts, not errors against biological measurements. The displayed digits preserve the recorded numerical values, without establishing biological accuracy.\cite{myohandNumericalRecord20261105}}
\label{tab:myohand-numerical-response}
\begin{tabular}{p{0.53\linewidth}p{0.26\linewidth}p{0.10\linewidth}}
\hline
Quantity & Contrast & Unit \\
\hline
Largest absolute final joint-coordinate separation, $\Delta q_{\infty}$ & $0.5027347840656727$ & rad \\
Selected final activation difference, $\Delta a$ & $0.05205729904534248$ & none \\
Selected final tendon-actuator force difference, $\Delta F$ & $-6.407722136774943$ & N \\
\hline
\end{tabular}
\end{table}


An independent standard-library audit of the recorded evidence completed 132 checks without a failed check and without rerunning the model. The required recorded numerical channels were finite, and recorded native warning counts were zero. The four stored endpoint mass-matrix spectra contained positive values. The raw mass matrices and residual operands were not independently exported, so the audit checked those stored diagnostics rather than solving their underlying matrices again. No contact was observed in this case. At the final control endpoint, the recorded soft-limit excess was $0.0004121829780273552\,\mathrm{rad}$; these soft constraints must not be described as exact hard joint bounds. The diagnostic gives no conservation proof, convergence study or physiological magnitude threshold.

The nonzero contrasts establish an input-to-state response in this configuration. They are differences between two simulated trajectories, not errors against human measurements. The 200 dependent integration steps are neither 200 subjects nor 200 independent experiments; no statistical $p$-value is inferred. Extra decimal places preserve the numerical record without establishing biological accuracy or repeatability across cases.

\subsubsection{The next claim needs a different reference}

The run used a CPU dynamics route. Readback showed a one-thread BLAS pool and a zero-sized dynamics pool, while transient compiler parallelism remained unmeasured. The 263 observed loaded-file identities form a bounded record, not a complete system closure. Resource and performance qualification remain unresolved: process high-water RSS was reported as 5,187,464 KiB and can retain pre-exec history, so it cannot be assigned to the hand model or used to establish a model-only 3 GiB bound. Tool/native timing near 448 ms and in-probe timing near 491.927 ms have unreconciled boundaries. Neither model RAM use, complete CPU/GPU partitioning, useful end-to-end latency nor a speedup follows from these values.

This mechanism supplies simulated activation, force and motion; it has no sensory measurements. Myoception---the programme's proposed sensory side of muscle-driven action---would require its own task-relevant sensory evidence. No afferent recording, learned controller, biological feedback loop, trained Large Movement Model or transfer result is demonstrated. The four proposed training-source groups remain neural observations, myography, measured motion and physics-generated examples. Before this generator could contribute training data, its relevant mechanics, controls, calibration, uncertainty, source rights and training lineage would need qualification, followed by independent held-out real tests. The amount/diversity hypothesis remains open. No protected test or quarantined source is reopened by this diagnostic.

The current record supplies no successful scaling curve or trained body-wide world model. LINK's eighteen fits and 720 saved epochs still fail both model-admission gates and all five benefit claims. New data would have to improve a new, properly qualified comparison; it cannot revise that negative result after the fact. The proposed programme becomes stronger when it names the particular relationship that survives a harder test, and equally when it identifies the one that fails.

% Prospective task/data protocol; no new experiment or admitted model is reported.
\section{From contact to coordinated action}

A ball leaves one hand and reaches the other. A prediction can follow both hands
and still miss the release, the contact or the order of events. Those differences
give the proposed paired study its next questions. We want an intent-to-action
mapping as close to one-to-one as the declared task permits; measuring the
intermediary gives us a way to test whether the ambiguity actually falls.

The proposed sequence moves from contact and comparatively simple movements to
bilateral object handling, punch discrimination and toy archery. It extends the
single-hand precision example without treating every task as another copy of a
reach. Each task needs its own anatomy, contact conditions, permissible context,
independent reference and error specification. The full snap remains a coordinated
arm, forearm, wrist and hand movement, including its preparation and return;
species-specific reference evidence must precede any claim about a monkey's
ability to perform it.

\subsection{Record one event across the instruments}

The proposed apparatus pairs scalp EEG with two surface-EMG recording sleeves,
one on each arm, and synchronized, calibrated left and right tripod views of the
hands and arms. OpenBCI or an equivalent configurable system remains a candidate.
Channel count, electrode sites and references, bandwidth, units, device delays
and clocks must be qualified for the chosen task before the recordings can enter
training. A nominal sampling rate cannot establish that the instruments observed
the same event at the same time.

These left/right views are the proposed hand-task configuration of the
motion-capture relay-composite. The earlier opposing top/bottom dot-map
configuration describes a different proposed apparatus. Each configuration
requires its own camera geometry, calibration, timing, coordinate frame,
visibility and occlusion checks. Joining two maps preserves their source and
uncertainty; it does not turn a model-filled gap into an observed point.

The musculoskeletal-simulation physics-\allowbreak{}engine would predict the body's response\linebreak[4]
within its qualified anatomy and mechanics. Independent movement observations
would test that response. Contact labels need an appropriate independent
observation, and a force claim needs a force reference. Surface EMG measures
muscle electrical activity: it supplies neither direct force nor sensory-nerve
activity, and it is not automatically the simulator's activation variable.

\subsection{Four tasks, four reasons to measure}

\paragraph{Touch and movement.}
Begin with rapid hand rubbing, then record the requested knuckle movement and
interlaced-hand stretch as separately identified conditions. The stretch extends
the arms with the palms turned outward. Rubbing combines contact and movement;
the comparison therefore needs matched conditions that can distinguish their
contributions, rather than labelling every change in EEG as touch. The
cold-hands reference supplies the intuitive action, not a cold-exposure result.

\paragraph{Transfer and juggling.}
Pass one juggling ball between the hands in the proposed hot-potato task, then
progress to the preferred two-ball juggling task. Release, flight and catch
connect the two limbs through the same object. Preserve the observed event
order, each limb's trajectory and the limits of contact observation. A predictor
that knows the task cue or the previous trajectory may succeed without useful
neural information, so those simpler explanations remain in the comparison.

\paragraph{Punch discrimination and NeuroBoxer.}
Define punch classes from independent movement observations, then test whether
the permitted recordings predict the class, trajectory and timing specified by
the protocol. Boxing also involves hand and wrist stabilization. NeuroBoxer is
the proposed virtual game built around these comparisons; a prime-era Mike Tyson
opponent is a proposed virtual game character. It supplies no physical opponent
or evidence of safety. Expert footage can inform movement and tactics within its rights and
camera limits. It supplies none of the expert's EEG, EMG or muscle-force data.

\paragraph{Toy archery and NeuroArcher.}
The proposed toy task uses arrows with small, cup-shaped suction tips and an
over-arm pouch with multiple arrows.
Grasp, draw, release and return give the two hands different roles under opposing
loads; finger and wrist changes matter as well as arm trajectories. NeuroArcher
would use that coordination to act on a virtual target. Neither the target
design nor a successful game action establishes measured biological mechanics.

\subsection{Test the neural contribution before assigning it control}

The first game milestone uses measured body motion to control the virtual body.
A later EEG-only milestone would test a genuine contribution from the permitted
brain observations. Its predictor must not receive withheld EMG or movement
targets. Eye activity, covert muscle activity, residual movement, shared cues,
trial pairing and future information need separately declared controls;
apparently hands-free play cannot settle those questions by appearance.

Keep the original brain-only, body-only and paired observation comparisons,
target-persistence controls and benefit criteria. Record permissible history and
causal windows before fitting preprocessing or choosing a model. Reserve the
required subject, session and task partitions before window construction; fit
transforms on training data only. Pretraining exposure must be resolved before
calling a final recording independent. Each reverse test predicts specified
recorded EEG features conditionally; it cannot identify a unique hidden brain
state or deliver sensation.

Qualified pretraining, paired fitting and fine-tuning would follow admitted data,
rights, clock, unit and baseline checks. Optional control-policy reinforcement
learning also needs declared states, actions, transitions and rewards with
appropriate coverage. A simulation reward tests the declared computational
objective; independent biological outcomes require their own evidence. Final
held-out observations stay outside rewards, tuning and model selection, and
frozen inference remains separate from independently evaluated adaptive updates.

\subsection{Separate the later writer from the recording proposal}

EEG and surface EMG record. Neuromuscular electrical stimulation (NMES) and functional electrical stimulation (FES) denote stimulation approaches that can
produce muscle contractions or functional movement; those are separate physical
output interfaces. Future sensation delivery likewise needs a qualified writer
or sensory transducer. Predicting a neural feature, commanding a virtual body
and producing a biological response answer different questions. Physical spinal
scope remains read-only.

The proposed later blindfolded condition changes visual feedback. It does not
independently validate neural control or retained learning. A retention test must
remove trainer information while keeping the necessary movement bridge and
comparable task conditions, then compare performance with the original controls.
No participant recording, stimulation, actuation or such retention experiment is
reported here. Physical writers and actuators remain outside the existing
acquisition/virtual-work budget scope. The retained LINK failures remain the
current computational result; this sequence specifies prospective comparisons.

% New prospective section. Not assembled, rendered, published or scientifically accepted.
% Preserve earlier studies, result tables, citations and frozen failures as historical evidence.
\section{Measure both ends. Test the translation.}

The following prospective protocol tests the paired learner described in ``Two pairs. One conditional learner'': can permitted observations on one side of a sensorimotor task predict withheld measurements on the other, on recordings the model has not seen? The apparatus remains non-invasive EEG, non-invasive surface EMG, a calibrated motion-capture relay-composite, an open musculoskeletal-simulation physics-engine with a lifelike explanatory view, and an open multimodal model. Its acquisition and validation are unfinished.

The earlier brain--muscle and measured-movement--virtual-body comparisons define the two pairs. Here the protocol specifies their records, targets, calibration and independent tests. Anatomical form constrains possible mechanisms; measured function must test the proposed connections. A conditional learner's combined prediction must answer to measurements outside itself.

\subsection{Keep the observations distinct}

Keep raw EEG, raw EMG, derived features and movement estimates in separate
records. Identify the participant and session in each record, then declare the
sensor sites, references, coordinate conventions, units, sampling and timestamp
semantics, and clocks. Trace every processing step and preserve missing intervals and uncertainty. Retain task
instructions, stimulus events, observed behaviour and any separately obtained
participant report. Independent participants or recordings cannot supply a
simultaneous brain--muscle--movement observation.

Surface EMG records muscle electrical activity. A sensory-loop claim also needs
known stimuli and appropriate measured feedback, contact or force channels;
muscle recordings do not provide incoming sensory-nerve traffic. Preserve the
actual spatial resolution of EEG measurements and source estimates. Templates
and explanatory cutaways do not register participants or localize individual
neurons.

\subsection{The spine shapes the loop}

The spinal cord changes how sensory input becomes muscle output. Account for
local coordination, sensory gating, reflex gain and some lasting adaptations,
rather than reducing it to a transmission delay. Separate a response to the
current posture or task from a change retained after training. Keep both apart
from adaptation in sensory receptors, peripheral pathways, neuromuscular
junctions and mechanical tissues. These processes interact; changed output alone
does not locate a memory.

Human H-reflex conditioning and the constrained mouse adaptation studied by
Lavaud and colleagues ask different questions.\cite{thompson2009retained,thompson2013retained,lavaud2024}
The human experiments do not establish brain-independent skill learning. The
mouse study does not establish unrestricted walking or reaching. Circuits
required for acquisition or recall are not thereby exclusive memory-storage
sites. Our human comparison uses the abstracts, selected results and limited
methods passages. The mouse comparison also draws on relevant author-manuscript
passages; equivalence to the final published edition and the complete methods
and statistics have not been established here. Passive proprioceptive training
motivates a peripheral observation model; its psychophysical and spindle-afferent
experiments were separate.\cite{ackerley2022retained} Those observations do not
establish complex memory stored in peripheral axons.

Define the function before choosing its computational substitute. Candidate
functions include gain and inhibition, phase- and load-dependent sensory
responses, coordination across supported muscles or joints, relevant rhythms,
and acquisition, retention and expression of a declared adaptation. Give each
function its own task, species, anatomical scope, measurement and reference.
Naming a mechanism does not show that the model implements it. A hand task
cannot inherit an unmeasured locomotor circuit.

Record task context and history. Distinguish the state within a trial from a
learned parameter change; retain elapsed time, intervening activity, resets,
drift and uncertainty. Recorded spinal or peripheral activity needs supported
channels of its own. Where those observations are absent, label internal state
as estimated or as a simulated controller. EEG, EMG and movement cannot be
relabeled as recorded spinal activity or identified interneurons.

Declare controls and endpoints before fitting or experimentation. Test
immediate context effects separately from retention after a declared interval
and from peripheral or tissue changes. Preserve the original held-out subjects
and tasks, eligible independent measurements, matched baselines and ablations.
Check alternative explanations, time shifts, missing inputs and artifact
leakage. Lower movement error alone cannot locate a biological adaptation,
establish its cause or certify useful plasticity.

To externalize a specified function, propose a substitute that can be tested
through declared peripheral and cerebral interfaces. Compare latency, state
dependence, stability, uncertainty and measured outcomes with the specified
biological function. Computational feedback, peripheral stimulation and restored
biological function require different evidence. A matched virtual trajectory
does not demonstrate spinal replacement, relocated learning or clinical repair.
Keep the spinal-cord site, measured function and proposed computational
substitute distinct. Physical spinal scope remains read-only.

\subsection{Give the two directions different targets}

For brain-to-limb prediction, give the model EEG and only the declared context
available at inference. Test its predicted muscle and movement observables
against withheld measurements. A brain-only condition receives no measured EMG
or movement input. If those inputs help, evaluate them in a separately labeled
hybrid condition.

For limb-to-brain prediction, give the model muscle activity, movement and
declared context. Its answer must allow for different patterns in the supported
recorded EEG features. A conditional distribution assigns probabilities to
specified feature values or ranges, given those inputs. Compare that prediction
with independently retained EEG, then report uncertainty and alternatives.
Movement can narrow the possibilities without identifying a unique hidden
brain state or subjective experience. Agreement between learned models is a
consistency check; it cannot supply the measured target.

\subsection{Reading a signal and writing an output}

The protocol records EEG and EMG; its two prediction directions do not themselves deliver sensation or enact a movement. ``Write sensation. Enact action'' defines the computer and physiological boundaries. Testing a writer would additionally require a separately qualified sensory writer or transducer, a specified actual response and distinct device, measurement, safety and clinical evidence. Indirect sensory feedback and direct neural writing remain different interfaces; a predicted EEG feature is neither. Enaction needs an appropriate virtual, robotic or biological output interface. A robot command has an action role, not the anatomical role of an efferent nerve.

Physical sensory writers, transducers and body actuators remain excluded from the acquisition and virtual-work plan, unqualified and unpriced. This proposal authorizes no stimulation, participant acquisition, biological actuation, new equipment purchase or training.

\subsection{Align before learning}

Calibrate the acquisition views. Verify their clocks against recorded events,
offsets and drift, while preserving occlusion and dropout. Keep a missing
observation missing: a predicted finger position is not a measurement. Declare
frames, scale, joint conventions and acceptable spatial, temporal and end-to-end
errors before selecting equipment or models. Position, joint, contact and force
claims each need an appropriate independent reference.

Specify how the virtual body represents subject scale, muscles, tendons, joint
constraints and contact. Declare its integration settings and observation
models. Biomechanical simulation predicts mechanics; VFX rendering makes those
predictions visible. Judge individual parts and their assembly against the
original anatomical, physical, functional and real-reference-series criteria.

\subsection{Build an open corpus without inventing pairs}

Inventory neural structure, neural activity, muscle structure and muscle
activity as separate domains. Use measured movement and task context as
constraints. For each source, record rights, version, access conditions,
species, participants, tasks, montages, units, clocks, preprocessing and missing
paired modalities. Broad open-data coverage remains a continuing objective,
not an assertion that every dataset is accessible, licensed, compatible or
already used. Check code, pretrained weights and training-data rights separately.

The candidate uses an encoder for each modality, explicit input-availability
masks, a shared conditional representation and separate decoders for the two
directions. Keep simple EEG-only, EMG-only, movement-only and hybrid baselines.
Define pretraining, post-training and fine-tuning tasks before fitting. An
existing encoder supplies a candidate component, not an already established
brain--body model.

\subsection{Make the reward answer to measurements}

First test paired supervised prediction in each direction. Compare any later\linebreak[4]
reinforcement-learning or other closed-loop refinement with that baseline.
Evaluate independent movement measurements, task outcomes, relevant contact or
force, latency and physical validity. Freeze reward units, terms, weights and
decision thresholds before evaluation. The candidate cannot invent the target
against which it earns a reward.

Backpropagation updates model parameters. Reverse prediction estimates another
observable. Measured biological feedback comes from the body. Identify which
process the study actually implements. A supervised prediction error alone does
not establish reinforcement learning: a proposed control study must also declare
its state, action, transition and reward.

Partition eligible participants, sessions and novel tasks before preprocessing
or model selection. Preserve TRAIN-only normalization, frozen seeds and
checkpoint selection, untouched held-out evaluation, uncertainty and original
negative controls. Include informative missing-modality, time-shift, artifact,
drift and alternative-inverse comparisons. Keep reserved recordings closed and
retain completed fits and failed benefit conjunctions. Evaluate frozen inference
separately from any proposed adaptive update. This prospective protocol admits
no new fitting, optimizer update or model/backend promotion.

\subsection{Show what was measured and what was inferred}

Show the selected source clock and the inputs available to each prediction.
Place recorded EEG, muscle activity and movement beside their predictions, with
uncertainty, errors, gaps and stimulus or contact context. Identify recordings,
inferred features and simulation as separate layers. Preserve site/channel
inspection and source-time play, pause, seek and step. Unsupported peripheral
or spinal pathways remain visible gaps. An animated cord cannot substitute for
a spinal recording or supported anatomical localization.

\section{Results and Remaining Questions}

The LINK comparison shows why an improved mean cannot stand in for a better decoder. Neither model meets the frozen requirements, and all five benefit claims fail. The earlier studies retain their narrower results: specified within-recording estimates, a separately calibrated different-day specialist, bounded simulated muscle correction and traceable geometry checks. They answer different questions and cannot reverse the LINK outcome.

The hand studies make the difference visible. The classifier identifies 116/150 attempted-finger labels. In the linked simulation, 54/54 issued commands are executed, yet only 41/60 source intentions become correct contacts and 13 unintended presses remain. A separate goal-driven policy completes 32/32 constrained simulated phrases without recorded brain activity driving that skill. Executing a decision, making the right decision and learning a phrase are different achievements. None establishes biological teaching, a neural policy for the phrases, voluntary whole-chain control, retained mastery or universal translation.

The next paired comparison would ask whether permitted neural observations predict withheld muscle or movement measurements better than the specified simpler controls. Unknown units, upstream clocks, contact identities, participant registration, joint and force measurements and missing nerve or spinal observations must be resolved or left explicitly unavailable. More computation cannot supply an absent measurement by assertion. Source support, independent implementation, final-checkpoint/update verification and reserved confirmation remain separate requirements; completing one cannot stand in for the others.

That experiment must keep the original physical, statistical, held-out, uncertainty and failure thresholds, with independent acceptance for the claimed domain and real-reference series for the applicable anatomy and function. Biological function and a complete working GUI require their own evidence. Publication and edition-review limits are collected in ``Evidence, scope and edition status''. The test returns to the opening distinction: does measuring the intermediary improve the declared intent-to-action mapping? If it yields useful assistance, does any learning remain after trainer information is removed under controlled conditions? Neither result is established here.

